% =============================================================================
%  FlowForge — Rapport du Mini-Projet Microservices
%  Module : Web Avancé & Microservices
%
%  COMPILATION : pdfLaTeX (2 passes pour la table des matières).
%  Sur Overleaf : téléverser tout le dossier `rapport/`, choisir report.tex
%  comme document principal, compilateur pdfLaTeX. Aucun paquet exotique,
%  aucun shell-escape.
%
%  Les diagrammes sont écrits en Mermaid (diagrams/*.mmd) puis pré-rendus en
%  PDF vectoriel par ./build-diagrams.sh. Le code source Mermaid de chaque
%  figure est reproduit en annexe C.
% =============================================================================
\documentclass[11pt,a4paper,oneside]{report}

% ---------- Encodage & langue -------------------------------------------------
\usepackage[utf8]{inputenc}
\usepackage[T1]{fontenc}
\usepackage[french]{babel}
\usepackage{lmodern}
\usepackage{textcomp}

% ---------- Mise en page ------------------------------------------------------
\usepackage[a4paper,top=2.5cm,bottom=2.5cm,left=2.5cm,right=2.5cm]{geometry}
\usepackage{fancyhdr}
\usepackage{titlesec}
\usepackage{setspace}
\onehalfspacing

% ---------- Graphiques, couleurs, tableaux ------------------------------------
\usepackage{graphicx}
\usepackage{xcolor}
\usepackage{booktabs}
\usepackage{longtable}
\usepackage{tabularx}
\usepackage{array}
\usepackage{float}
\usepackage{caption}
\usepackage{enumitem}
\usepackage{tcolorbox}
\tcbuselibrary{skins,breakable}

% ---------- Code source -------------------------------------------------------
\usepackage{listings}

% ---------- Liens (à charger en dernier) --------------------------------------
\usepackage[hidelinks,breaklinks=true]{hyperref}

% =============================================================================
%  >>> À REMPLIR PAR L'ÉQUIPE — ce sont les SEULES lignes à modifier <<<
% =============================================================================
\newcommand{\membreUn}{AMAZIGH Hicham}
\newcommand{\membreDeux}{BOUDILI Mehdi}
\newcommand{\membreTrois}{CHIBOUB Abdelaziz}
\newcommand{\membreQuatre}{SENDID Ahmed}
\newcommand{\filiere}{Master Informatique, Gouvernance et Transformation Digitale}
\newcommand{\anneeUniv}{2025 -- 2026}
\newcommand{\groupe}{Groupe 15}
\newcommand{\encadrant}{Pr. Hamza EL KHATEB}
\newcommand{\etablissement}{Université Mohammed V --- Faculté des Sciences de Rabat}
\newcommand{\moduleTitre}{Microservices}
% =============================================================================

\newcommand{\projet}{FlowForge}
\newcommand{\soustitre}{Plateforme collaborative de gestion de tâches}

% ---------- Palette -----------------------------------------------------------
\definecolor{ffblue}{HTML}{1D4ED8}
\definecolor{ffdark}{HTML}{1E3A8A}
\definecolor{ffgray}{HTML}{64748B}
\definecolor{ffbg}{HTML}{F1F5F9}
\definecolor{ffgreen}{HTML}{15803D}
\definecolor{fforange}{HTML}{B45309}
\definecolor{ffred}{HTML}{B91C1C}
\definecolor{codebg}{HTML}{F8FAFC}
\definecolor{codekw}{HTML}{1D4ED8}
\definecolor{codestr}{HTML}{15803D}
\definecolor{codecom}{HTML}{64748B}

% ---------- En-têtes / pieds de page -----------------------------------------
\pagestyle{fancy}
\fancyhf{}
% En-tête : nom du projet à gauche, titre du chapitre courant à droite.
% Le sous-titre a été retiré de l'en-tête : à gauche « FlowForge — Plateforme
% collaborative... » et à droite le titre de chapitre étaient tous deux longs et
% se chevauchaient au milieu de la ligne. \nouppercase évite l'affichage en
% capitales du \leftmark, plus compact et lisible.
\fancyhead[L]{\small\textcolor{ffgray}{\projet}}
\fancyhead[R]{\small\textcolor{ffgray}{\nouppercase{\leftmark}}}
\fancyfoot[C]{\thepage}
\renewcommand{\headrulewidth}{0.4pt}
\renewcommand{\footrulewidth}{0pt}
\fancypagestyle{plain}{\fancyhf{}\fancyfoot[C]{\thepage}\renewcommand{\headrulewidth}{0pt}}

% ---------- Style des titres --------------------------------------------------
\titleformat{\chapter}[display]
  {\normalfont\huge\bfseries\color{ffdark}}
  {\normalsize\color{ffblue}\MakeUppercase{\chaptertitlename}\ \thechapter}
  {8pt}{\Huge}
\titlespacing*{\chapter}{0pt}{-20pt}{28pt}
\titleformat{\section}{\normalfont\Large\bfseries\color{ffdark}}{\thesection}{1em}{}
\titleformat{\subsection}{\normalfont\large\bfseries\color{ffblue}}{\thesubsection}{1em}{}

% =============================================================================
%  Styles de code (listings)
% =============================================================================
\lstdefinelanguage{TypeScript}{
  keywords={abstract,as,async,await,break,case,catch,class,const,constructor,continue,
    declare,default,delete,do,else,enum,export,extends,false,finally,for,from,function,
    get,if,implements,import,in,instanceof,interface,let,new,null,private,protected,
    public,readonly,return,set,static,super,switch,this,throw,true,try,type,typeof,
    undefined,var,void,while,yield},
  keywordstyle=\color{codekw}\bfseries,
  ndkeywords={string,number,boolean,any,Promise,Record,Date,Buffer},
  ndkeywordstyle=\color{fforange},
  sensitive=true,
  comment=[l]{//},
  morecomment=[s]{/*}{*/},
  commentstyle=\color{codecom}\itshape,
  stringstyle=\color{codestr},
  morestring=[b]',
  morestring=[b]",
  morestring=[b]`,
}

\lstdefinelanguage{yaml}{
  keywords={true,false,null,yes,no,on,off},
  keywordstyle=\color{fforange},
  sensitive=false,
  comment=[l]{\#},
  commentstyle=\color{codecom}\itshape,
  morestring=[b]',
  morestring=[b]",
  stringstyle=\color{codestr},
}

% `listings` ne connaît pas JSON nativement : on le déclare.
\lstdefinelanguage{json}{
  keywords={true,false,null},
  keywordstyle=\color{fforange},
  sensitive=false,
  comment=[l]{//},
  commentstyle=\color{codecom}\itshape,
  morestring=[b]",
  stringstyle=\color{codestr},
}

\lstdefinelanguage{nginxconf}{
  keywords={server,location,upstream,http,events,listen,proxy_pass,proxy_set_header,
    proxy_http_version,proxy_read_timeout,server_name,return,default_type,
    client_max_body_size},
  keywordstyle=\color{codekw}\bfseries,
  sensitive=true,
  comment=[l]{\#},
  commentstyle=\color{codecom}\itshape,
  morestring=[b]",
  stringstyle=\color{codestr},
}

\lstdefinestyle{ffcode}{
  backgroundcolor=\color{codebg},
  basicstyle=\ttfamily\scriptsize,
  keywordstyle=\color{codekw}\bfseries,
  commentstyle=\color{codecom}\itshape,
  stringstyle=\color{codestr},
  numbers=left,
  numberstyle=\tiny\color{ffgray},
  numbersep=8pt,
  frame=leftline,
  framerule=2pt,
  rulecolor=\color{ffblue},
  xleftmargin=14pt,
  breaklines=true,
  breakatwhitespace=false,
  showstringspaces=false,
  tabsize=2,
  columns=fullflexible,
  upquote=true,
  captionpos=b,
  extendedchars=true,
  literate=%
    {é}{{\'e}}1 {è}{{\`e}}1 {ê}{{\^e}}1 {ë}{{\"e}}1
    {à}{{\`a}}1 {â}{{\^a}}1 {ä}{{\"a}}1
    {î}{{\^i}}1 {ï}{{\"i}}1
    {ô}{{\^o}}1 {ö}{{\"o}}1
    {ù}{{\`u}}1 {û}{{\^u}}1 {ü}{{\"u}}1
    {ç}{{\c c}}1
    {É}{{\'E}}1 {È}{{\`E}}1 {Ê}{{\^E}}1
    {À}{{\`A}}1 {Ç}{{\c C}}1 {Ô}{{\^O}}1
    {«}{{\guillemotleft}}1 {»}{{\guillemotright}}1
    {°}{{\textdegree}}1 {€}{{\texteuro}}1
    {–}{{--}}1 {—}{{---}}1 {’}{{'}}1 {…}{{\ldots}}1
    % Symboles présents dans les sources Mermaid reproduites en annexe C
    {·}{{\textperiodcentered}}1 {•}{{\textbullet}}1
    {→}{{$\rightarrow$}}1 {←}{{$\leftarrow$}}1 {↓}{{$\downarrow$}}1
    {∈}{{$\in$}}1 {×}{{$\times$}}1 {≥}{{$\geq$}}1 {≤}{{$\leq$}}1
}
\lstset{style=ffcode}
\renewcommand{\lstlistingname}{Extrait}

% =============================================================================
%  Macros de présentation
% =============================================================================

% --- Encadré « point clé » ----------------------------------------------------
% Le bandeau de titre est rempli en BLEU FONCÉ : le texte du titre doit donc
% être BLANC. Le corps, lui, reste sur fond clair avec du texte foncé.
\newtcolorbox{pointcle}[1][]{
  enhanced, breakable,
  colback=ffbg, colframe=ffblue,
  colbacktitle=ffblue,          % fond du bandeau de titre
  coltitle=white,               % texte du titre : BLANC (lisible)
  fonttitle=\bfseries,
  coltext=ffdark,               % texte du corps : foncé sur fond clair
  boxrule=0pt, leftrule=3pt, arc=1pt,
  title={#1},
  left=8pt, right=8pt, top=6pt, bottom=6pt
}

% --- Encadré « choix d'architecture » ----------------------------------------
\newtcolorbox{choix}[1][]{
  enhanced, breakable,
  colback=white, colframe=ffgreen,
  colbacktitle=ffgreen,         % fond du bandeau de titre
  coltitle=white,               % texte du titre : BLANC (lisible)
  fonttitle=\bfseries,
  coltext=black,                % texte du corps : noir sur fond blanc
  boxrule=0pt, leftrule=3pt, arc=1pt,
  title={#1},
  left=8pt, right=8pt, top=6pt, bottom=6pt
}

% --- Figure Mermaid -----------------------------------------------------------
% #1 = nom du fichier (sans extension) dans diagrams/
% #2 = légende  |  #3 = étiquette  |  #4 = largeur relative
\newcommand{\mermaidfig}[4]{%
  \begin{figure}[H]
    \centering
    \includegraphics[width=#4\linewidth]{diagrams/#1.pdf}
    \caption{#2}
    \label{fig:#3}
  \end{figure}%
}

% --- Emplacement de capture d'écran ------------------------------------------
% Si le fichier screenshots/<nom>.png existe, il est inclus automatiquement.
% Sinon, un cadre « à insérer » est affiché à sa place.
\newcommand{\phbox}[1]{%
  \setlength{\fboxsep}{0pt}%
  \fcolorbox{ffblue}{ffbg}{%
    \begin{minipage}[c][52mm][c]{\linewidth}
      \centering
      {\bfseries\large\color{ffblue} CAPTURE D'ÉCRAN À INSÉRER}\\[3mm]
      {\ttfamily\small\color{ffdark} rapport/screenshots/#1.png}\\[3mm]
      {\small\itshape\color{ffgray} Déposer le fichier à cet emplacement :\\
       il sera inclus automatiquement à la prochaine compilation.}
    \end{minipage}%
  }%
}
\newcommand{\capture}[3]{%
  \begin{figure}[H]
    \centering
    \IfFileExists{screenshots/#1.png}
      {\fbox{\includegraphics[width=\linewidth]{screenshots/#1.png}}}
      {\phbox{#1}}
    \caption{#2}
    \label{fig:cap-#3}
  \end{figure}%
}

\newcommand{\code}[1]{\texttt{\small #1}}

% =============================================================================
\begin{document}
\sloppy

% =============================================================================
%  PAGE DE GARDE
% =============================================================================
\begin{titlepage}
\centering
% L'interligne 1,5 du corps du rapport gonflerait la page de garde et
% pousserait l'encadrant sur une 2e page : on repasse en interligne simple.
\singlespacing
\vspace*{0.2cm}

% Logo de la faculté. \IfFileExists évite de casser la compilation si le
% fichier venait à manquer : la page de garde se contenterait alors du texte.
\IfFileExists{screenshots/00-logo_fsr.png}
  {\includegraphics[width=5.2cm]{screenshots/00-logo_fsr.png}\par
   \vspace{0.45cm}}
  {}

{\large\scshape \etablissement\par}
\vspace{0.25cm}
{\large\scshape \filiere\par}
\vspace{0.15cm}
{\large Année universitaire \anneeUniv\par}

\vspace{0.85cm}

{\color{ffgray}\rule{\linewidth}{0.4pt}}
\vspace{0.4cm}

{\small\bfseries\color{ffblue}\MakeUppercase{Module — \moduleTitre}\par}
\vspace{0.35cm}

{\Huge\bfseries\color{ffdark} Mini-Projet Microservices\par}
\vspace{0.55cm}

{\huge\bfseries\color{ffblue} \projet\par}
\vspace{0.25cm}
{\Large\itshape \soustitre\par}

\vspace{0.4cm}
{\color{ffgray}\rule{\linewidth}{0.4pt}}

\vspace{0.9cm}

\begin{tcolorbox}[width=0.95\linewidth, colback=ffbg, colframe=ffblue,
                  boxrule=0.6pt, arc=2pt, left=8pt, right=8pt, top=8pt, bottom=8pt]
\centering
{\bfseries\color{ffdark} Architecture réalisée\par}
\vspace{2mm}
{\small
11 microservices conteneurisés \quad\textbullet\quad 3 technologies backend\\
3 SGBD (PostgreSQL, MySQL, Elasticsearch) \quad\textbullet\quad Bus d'événements RabbitMQ\\
Passerelle API NGINX \quad\textbullet\quad Frontend SPA Vue 3\par}
\end{tcolorbox}

\vfill

{\large\bfseries Réalisé par :\par}
\vspace{0.3cm}
{\large
\membreUn\\[1.5mm]
\membreDeux\\[1.5mm]
\membreTrois\\[1.5mm]
\membreQuatre\par}

\vspace{0.55cm}
{\large\bfseries \groupe\par}

\vspace{0.6cm}
{\large\bfseries Encadré par :\par}
\vspace{0.15cm}
{\large \encadrant\par}

\vspace{0.4cm}
\end{titlepage}

% =============================================================================
\tableofcontents
\listoffigures
\cleardoublepage

% =============================================================================
\chapter{Introduction}
% =============================================================================

\section{Contexte du projet}

Ce rapport présente \projet, une plateforme collaborative de gestion de tâches
développée dans le cadre du module \emph{Web Avancé \& Microservices}. Le sujet
demandait de concevoir une plateforme permettant à des utilisateurs de
s'inscrire, de gérer leurs tâches, de les organiser par catégories, de les
commenter et de consulter des statistiques de productivité, le tout accessible
via un frontend et une passerelle API centralisée.

Nous avons choisi de traiter ce sujet non pas comme un simple exercice de
découpage en six services, mais comme la construction d'un système réellement
distribué : chaque service possède sa propre base, communique avec les autres
soit par API HTTP soit par un bus d'événements asynchrone, et reste déployable
et redémarrable indépendamment.

\begin{pointcle}[Ce qui distingue ce projet]
Au-delà des six microservices exigés, la plateforme intègre un
\textbf{bus d'événements RabbitMQ} (exchange topic, files durables,
\emph{dead-letter queues} et réessais), un \textbf{moteur de recherche
Elasticsearch} alimenté par événements avec filtrage des droits poussé dans le
moteur, un \textbf{service d'emailing} idempotent, un service de
\textbf{pièces jointes} adossé à un stockage objet, et des
\textbf{notifications temps réel par WebSocket}. Ces éléments font l'objet du
chapitre~\ref{chap:extras}.
\end{pointcle}

\section{Objectifs}

\begin{itemize}[leftmargin=*]
  \item Concevoir une architecture microservices cohérente, où le découpage
        suit les \emph{frontières métier} et non des considérations techniques.
  \item Garantir l'\textbf{indépendance des données} : une base dédiée par
        service, aucune lecture croisée.
  \item Mettre en place une \textbf{authentification sans état} (JWT) validée
        localement par chaque service.
  \item Exposer un \textbf{point d'entrée unique} (NGINX) servant à la fois la
        SPA et le \emph{reverse proxy} de chaque service.
  \item Démontrer la maîtrise de \textbf{plusieurs écosystèmes} : Python,
        TypeScript et Java coexistent derrière un contrat d'API et
        d'événements commun.
  \item Assurer la \textbf{résilience} des traitements asynchrones (réessais,
        files de rebut, idempotence).
\end{itemize}

\section{Périmètre fonctionnel}

\begin{longtable}{@{}p{4.2cm}p{10.3cm}@{}}
\caption{Périmètre fonctionnel couvert}\label{tab:perimetre}\\
\toprule
\textbf{Domaine} & \textbf{Fonctionnalités couvertes} \\
\midrule
\endfirsthead
\toprule
\textbf{Domaine} & \textbf{Fonctionnalités couvertes} \\
\midrule
\endhead
\bottomrule
\endlastfoot
Authentification &
  Inscription, connexion, JWT (accès + rafraîchissement), déconnexion globale,
  réinitialisation de mot de passe, vérification d'email, double authentification
  (TOTP), gestion et révocation des sessions, verrouillage après échecs répétés. \\
\addlinespace
Profils &
  Consultation et mise à jour du profil, photo de profil (validation de type et
  de taille), biographie, fuseau horaire, langue, préférences libres. \\
\addlinespace
Tâches (TODO) &
  CRUD complet, priorité, statut, date d'échéance, sous-tâches, étiquettes,
  assignations multiples, dépendances entre tâches, récurrence, rappels,
  colonnes kanban, filtres enregistrés, journal d'activité. \\
\addlinespace
Catégories &
  CRUD, membres et rôles propres à la catégorie, invitations par jeton à usage
  unique, archivage, suppression logique. \\
\addlinespace
Commentaires &
  Ajout, consultation, modification et suppression logique des commentaires,
  avec contrôle de visibilité délégué au service des tâches. \\
\addlinespace
Statistiques &
  Taux de complétion, totaux, répartition par statut / catégorie / priorité,
  tendance de complétion sur les derniers jours, filtres combinables. \\
\addlinespace
Notifications \newline (\emph{bonus}) &
  Notifications in-app persistées, poussées en temps réel par WebSocket,
  détection des échéances proches (moins de 24~h) par tâche planifiée,
  compteur de non-lues, marquage lu / tout lu, suppression. \\
\end{longtable}

\section{Périmètre technique et correspondance avec le cahier des charges}

Le sujet imposait six microservices, plus un septième optionnel valorisé en
bonus. Le tableau~\ref{tab:mapping} établit la correspondance exacte entre les
services demandés et ceux que nous avons livrés.

\begin{table}[H]
\centering
\small
\begin{tabularx}{\linewidth}{@{}p{0.5cm}p{3.4cm}p{3.5cm}X@{}}
\toprule
\textbf{\#} & \textbf{Service demandé} & \textbf{Service livré} & \textbf{Technologie / base} \\
\midrule
1 & Auth-Service       & \code{auth-service}       & Django 5.2 + DRF \textbullet{} PostgreSQL \code{auth} \\
2 & User-Service       & \code{user-service}       & Django 5.2 + DRF \textbullet{} PostgreSQL \code{users} \\
3 & TODO-Service       & \code{task-service}       & NestJS 11 + TypeORM \textbullet{} PostgreSQL \code{tasks} \\
4 & Category-Service   & \code{category-service}   & NestJS 11 + TypeORM \textbullet{} PostgreSQL \code{categories} \\
5 & Comment-Service    & \code{comment-service}    & NestJS 11 + TypeORM \textbullet{} PostgreSQL \code{comments} \\
6 & Statistics-Service & \code{statistics-service} & NestJS 11 \textbullet{} \emph{sans état, aucune base} \\
\midrule
\textbf{B} & \textbf{Notification-Service} \newline \emph{(bonus)} & \code{notification-service} & Spring Boot 3.4 + WebSocket \textbullet{} MySQL \code{notifications} \\
\bottomrule
\end{tabularx}
\caption{Correspondance entre les services exigés et les services livrés}
\label{tab:mapping}
\end{table}

À ces sept services s'ajoutent quatre services supplémentaires, non demandés,
présentés au chapitre~\ref{chap:extras} : \code{workspace-service},
\code{attachment-service}, \code{email-service} et \code{search-service}.

\subsection{Respect des contraintes techniques obligatoires}

\begin{table}[H]
\centering
\small
\begin{tabularx}{\linewidth}{@{}p{5.6cm}p{1.1cm}X@{}}
\toprule
\textbf{Contrainte} & \textbf{État} & \textbf{Réalisation} \\
\midrule
Minimum 2 technologies différentes & \textbf{Oui} &
  \textbf{3} technologies : Python/Django, TypeScript/NestJS, Java/Spring Boot. \\
\addlinespace
Au moins un service par technologie & \textbf{Oui} &
  2 services Django, 4 services NestJS, 5 services Spring Boot. \\
\addlinespace
Minimum 2 SGBD différents & \textbf{Oui} &
  \textbf{3} SGBD : PostgreSQL 16, MySQL 8, Elasticsearch 8.15. \\
\addlinespace
Base dédiée par service & \textbf{Oui} &
  9 bases dans 9 conteneurs distincts. Aucun service ne lit la base d'un autre. \\
\addlinespace
Frontend obligatoire & \textbf{Oui} &
  SPA Vue 3 + TypeScript + Vite (framework libre). \\
\addlinespace
Frontend via la passerelle uniquement & \textbf{Oui} &
  Le conteneur \code{frontend} ne publie aucun port : il n'est joignable que
  par NGINX. \\
\addlinespace
Frontend servi depuis la racine \code{/} & \textbf{Oui} &
  \code{location / \{ proxy\_pass http://frontend\_app; \}} (annexe~\ref{ann:nginx}). \\
\addlinespace
Tous les services conteneurisés & \textbf{Oui} &
  19 conteneurs orchestrés par un unique \code{docker-compose.yml}. \\
\addlinespace
Reverse proxy pour chaque service & \textbf{Oui} &
  Un bloc \code{location} par microservice exposant une API HTTP.
  Seul l'\code{email-service} n'en a pas : il n'expose \emph{aucune} route
  publique et n'est piloté que par les événements RabbitMQ. Lui ouvrir une
  porte sur la passerelle aurait été une surface d'attaque gratuite. \\
\bottomrule
\end{tabularx}
\caption{Conformité aux contraintes techniques obligatoires}
\end{table}

\section{Architecture retenue}

Trois décisions structurent l'ensemble du système.

\begin{choix}[1. Une base de données par service, sans exception]
Chaque microservice possède sa propre base, dans son propre conteneur. Aucune
jointure SQL n'existe entre deux domaines : les références inter-services
(\code{task\_id}, \code{workspace\_id}, \code{author\_user\_id}\ldots) sont de
simples entiers, sans contrainte de clé étrangère, puisqu'ils désignent des
lignes situées dans une \emph{autre} base. Le coût est réel — il faut résoudre
ces liens par API ou par événement — mais le bénéfice l'est aussi : un service
peut migrer son schéma, voire changer de SGBD, sans coordination avec les
autres. C'est exactement ce qui nous a permis de basculer les services
\code{notification} et \code{email} sur MySQL sans toucher au reste.
\end{choix}

\begin{choix}[2. Authentification sans état par JWT à secret partagé]
L'\code{auth-service} est le seul émetteur de jetons. Il les signe en HS256
avec un secret (\code{JWT\_SIGNING\_KEY}) partagé par tous les services. Chaque
service valide donc le jeton \emph{localement}, sans appel réseau vers
l'auth-service : il n'y a ni goulot d'étranglement, ni point de défaillance
unique sur le chemin d'authentification. La revendication \code{user\_id}
suffit à identifier l'appelant.
\end{choix}

\begin{choix}[3. Communication synchrone pour les décisions, asynchrone pour les effets]
Quand un service a besoin d'une réponse \emph{avant} de pouvoir agir — par
exemple : « cet utilisateur a-t-il le droit de commenter cette tâche ? » — il
fait un appel HTTP direct et attend. En revanche, tout ce qui est un
\emph{effet de bord} — indexer la tâche dans Elasticsearch, notifier l'assigné,
envoyer un email de bienvenue — passe par RabbitMQ. Le producteur n'attend pas,
n'a pas besoin de savoir qui consomme, et l'ajout d'un nouveau consommateur ne
modifie aucune ligne de code chez lui.
\end{choix}

% =============================================================================
\chapter{Schéma de l'architecture globale}
\label{chap:archi}
% =============================================================================

\section{Vue d'ensemble}

La figure~\ref{fig:archi} présente l'architecture complète : le frontend, la
passerelle NGINX, les onze microservices, les trois SGBD et les flux de
communication (synchrones en trait plein, asynchrones en pointillés).

\mermaidfig{01-architecture-globale}{Architecture globale de la plateforme \projet{} — une base de données dédiée par service}{archi}{1.0}

Le navigateur ne connaît qu'une seule adresse : \code{http://localhost:8080}.
NGINX y sert la SPA depuis la racine et relaie chaque préfixe \code{/api/...}
vers le microservice compétent. Aucun service backend n'est joignable
directement depuis l'extérieur du réseau Docker.

\begin{pointcle}[Aucune base n'est partagée]
La figure~\ref{fig:archi} se lit ainsi : \textbf{à chaque service correspond
exactement une flèche, vers exactement une base}. Il n'existe nulle part une
base commune sur laquelle deux services viendraient se brancher. Sept services
possèdent leur propre instance PostgreSQL, deux possèdent leur propre instance
MySQL, le \code{search-service} possède son cluster Elasticsearch — et le
\code{statistics-service} n'a \emph{aucune} base, parce qu'il ne détient aucune
donnée.

Concrètement, cela signifie qu'aucun service ne peut lire ni écrire dans les
tables d'un autre : la seule façon d'obtenir une donnée qui ne vous appartient
pas est de la \emph{demander} (appel HTTP) ou de l'\emph{apprendre} (événement
RabbitMQ). C'est contraignant à écrire, mais c'est ce qui garantit qu'un
changement de schéma dans un service ne casse jamais un autre service.
\end{pointcle}

\section{Routage de la passerelle}

\mermaidfig{02-routage-nginx}{Table de routage de la passerelle NGINX}{routage}{1.0}

\begin{table}[H]
\centering
\small
\begin{tabularx}{\linewidth}{@{}p{3.9cm}p{3.9cm}X@{}}
\toprule
\textbf{Préfixe} & \textbf{Service cible} & \textbf{Rôle} \\
\midrule
\code{/} & \code{frontend:9000} & SPA Vue 3 (fallback \code{index.html}) \\
\code{/api/auth/} & \code{auth-service} & Inscription, connexion, JWT, sessions \\
\code{/api/profiles/} & \code{user-service} & Profils utilisateurs \\
\code{/api/tasks/} & \code{task-service} & Tâches et métadonnées \\
\code{/api/categories/} & \code{category-service} & Catégories, membres, invitations \\
\code{/api/comments/} & \code{comment-service} & Commentaires \\
\code{/api/statistics/} & \code{statistics-service} & Métriques de productivité \\
\code{/api/workspaces/} & \code{workspace-service} & Espaces de travail \\
\code{/api/attachments/} & \code{attachment-service} & Pièces jointes \\
\code{/api/notifications/} & \code{notification-service} & Notifications in-app \\
\code{/ws/notifications} & \code{notification-service} & Socket temps réel (\emph{Upgrade}) \\
\code{/api/search/} & \code{search-service} & Recherche globale \\
\bottomrule
\end{tabularx}
\caption{Routage par préfixe d'URL}
\end{table}

\begin{pointcle}[Le cas particulier du WebSocket]
NGINX ne relaie pas par défaut l'en-tête \code{Upgrade}. Sans les deux lignes
\code{proxy\_set\_header Upgrade \$http\_upgrade;} et
\code{proxy\_set\_header Connection "upgrade";}, la négociation WebSocket
échoue silencieusement. Le \code{proxy\_read\_timeout} est par ailleurs porté à
une heure, faute de quoi NGINX fermerait la socket après 60~secondes
d'inactivité — alors qu'une session de notification reste ouverte tant que
l'onglet l'est.
\end{pointcle}

\section{Topologie des conteneurs}

\mermaidfig{03-topologie-docker}{Topologie Docker Compose : conteneurs, images, volumes}{docker}{1.0}

L'ensemble tient dans un seul réseau \emph{bridge}
(\code{flowforge-platform-network}). Les conteneurs se résolvent entre eux par
leur nom de service grâce au DNS embarqué de Docker : \code{task-service}
contacte \code{http://category-service:8000} sans jamais connaître d'adresse IP.
Chaque base est adossée à un volume nommé, si bien que les données survivent à
un \code{docker compose down}.

\section{Flux de communication}

\subsection{Inscription : une chaîne de provisionnement asynchrone}

L'inscription illustre le mieux le fonctionnement du bus. L'\code{auth-service}
se contente de créer le compte et de publier un unique événement
\code{user.registered}. Trois services réagissent \emph{indépendamment} : le
profil est créé, l'espace de travail personnel est provisionné, l'email de
bienvenue est envoyé.

\mermaidfig{04-sequence-inscription}{Séquence d'inscription et de provisionnement asynchrone}{seq-inscription}{0.95}

\begin{pointcle}[Publication après validation de la transaction]
L'événement est publié dans un \code{transaction.on\_commit()} : si
l'insertion en base échoue et que la transaction est annulée, \emph{aucun}
événement n'est émis. On ne notifie jamais la création d'un utilisateur qui
n'existe pas. C'est une garantie simple mais essentielle à la cohérence du
système.
\end{pointcle}

\subsection{Cycle de vie d'une tâche}

\mermaidfig{05-sequence-tache}{Création, indexation et assignation d'une tâche}{seq-tache}{0.95}

\subsection{Bus d'événements}

\mermaidfig{06-bus-evenements}{Bus d'événements RabbitMQ : producteurs, exchange topic, files et rebut}{bus}{1.0}

% =============================================================================
\chapter{Code source des backends}
\label{chap:backends}
% =============================================================================

Ce chapitre documente les sept services du périmètre exigé (les six services
demandés et le service de notification, en bonus). Pour chacun : la pile
technique, la structure du projet, le modèle de données, les points d'entrée
exposés et un extrait de code commenté représentatif.

\section{Vue d'ensemble des piles techniques}

\begin{table}[H]
\centering
\small
\begin{tabularx}{\linewidth}{@{}p{3.4cm}p{3.3cm}p{2.4cm}X@{}}
\toprule
\textbf{Service} & \textbf{Pile} & \textbf{Port} & \textbf{Persistance} \\
\midrule
\code{auth-service}         & Django 5.2 + DRF       & 8002 & PostgreSQL \code{auth} \\
\code{user-service}         & Django 5.2 + DRF       & 8001 & PostgreSQL \code{users} \\
\code{task-service}         & NestJS 11 + TypeORM    & 8003 & PostgreSQL \code{tasks} \\
\code{category-service}     & NestJS 11 + TypeORM    & 8009 & PostgreSQL \code{categories} \\
\code{comment-service}      & NestJS 11 + TypeORM    & 8010 & PostgreSQL \code{comments} \\
\code{statistics-service}   & NestJS 11              & 8011 & \emph{aucune} \\
\code{notification-service} & Spring Boot 3.4 / Java 21 & 8006 & MySQL \code{notifications} \\
\bottomrule
\end{tabularx}
\caption{Piles techniques des services du périmètre exigé}
\end{table}

% -----------------------------------------------------------------------------
\section{Auth-Service}
% -----------------------------------------------------------------------------

\subsection{Rôle}

L'\code{auth-service} est la racine de confiance de la plateforme. Il possède
les comptes, les mots de passe et les sessions ; il est le \emph{seul} émetteur
de JWT. Aucun autre service ne stocke de mot de passe, et aucun autre service
n'a besoin de l'appeler pour valider un jeton.

\subsection{Pile technique}

Django~5.2 avec Django REST Framework et \emph{SimpleJWT}. Le choix de Django se
justifie ici par son système d'authentification éprouvé : hachage PBKDF2,
gestion des permissions et surtout un modèle utilisateur extensible, que nous
avons enrichi pour la 2FA et le verrouillage de compte.

\subsection{Structure du projet}

\begin{lstlisting}[language=bash,numbers=none]
server/auth-service/
  auth_django/            # configuration du projet Django
    settings.py           #   BD, SimpleJWT, ALLOWED_HOSTS
    cors.py               #   middleware CORS ecrit a la main
  auth_app/
    models.py             #   User, AuthSession, AuthToken, LoginAttempt
    serializers.py        #   validation des entrees (DRF)
    services.py           #   logique metier : login, 2FA, jetons
    views.py              #   points d'entree HTTP
    events.py             #   publication RabbitMQ
    migrations/
  Dockerfile
\end{lstlisting}

\subsection{Modèle de données}

\mermaidfig{08-modele-auth-user}{Modèle de données des services Auth et User}{modele-auth}{0.85}

Quatre tables. \code{auth\_app\_user} étend \code{AbstractUser} de Django et
ajoute la vérification d'email, le compteur d'échecs, le verrouillage temporaire
et les deux colonnes de 2FA. \code{auth\_app\_authsession} matérialise chaque
jeton de rafraîchissement émis (ce qui permet de lister et révoquer les sessions
une à une). \code{auth\_app\_authtoken} porte, dans une table unique
discriminée par \code{purpose}, les jetons de vérification d'email, de
réinitialisation de mot de passe et d'activation. \code{auth\_app\_loginattempt}
journalise chaque tentative de connexion.

\begin{pointcle}[Les jetons ne sont jamais stockés en clair]
\code{AuthToken} ne conserve que le \textbf{haché SHA-256} du jeton
(\code{token\_hash}, unique). Le jeton brut n'est retourné qu'une seule fois, à
l'émission. Une fuite de la base ne permettrait donc pas de réinitialiser le mot
de passe d'un utilisateur. Le même principe est appliqué aux invitations
(catégories et espaces de travail).
\end{pointcle}

\subsection{Points d'entrée exposés}

\begin{longtable}{@{}p{1.4cm}p{6.1cm}p{1.3cm}p{5.1cm}@{}}
\caption{Points d'entrée de l'Auth-Service}\label{tab:ep-auth}\\
\toprule
\textbf{Méth.} & \textbf{Chemin} & \textbf{Auth} & \textbf{Rôle} \\
\midrule
\endfirsthead
\toprule
\textbf{Méth.} & \textbf{Chemin} & \textbf{Auth} & \textbf{Rôle} \\
\midrule
\endhead
\bottomrule
\endlastfoot
POST & \code{/api/auth/register/} & Non & Crée le compte, publie \code{user.registered} \\
POST & \code{/api/auth/login/} & Non & Retourne \code{\{access, refresh\}} ou exige un code 2FA \\
POST & \code{/api/auth/refresh/} & Non & Renouvelle le jeton d'accès \\
POST & \code{/api/auth/logout/} & Oui & Révoque une session \\
POST & \code{/api/auth/logout-all/} & Oui & Révoque toutes les sessions \\
POST & \code{/api/auth/password/change/} & Oui & Change le mot de passe et déconnecte partout \\
POST & \code{/api/auth/password/forgot/} & Non & Émet un jeton de réinitialisation \\
POST & \code{/api/auth/password/reset/confirm/} & Non & Applique le nouveau mot de passe \\
POST & \code{/api/auth/email/verify/request/} & Oui & Émet un jeton de vérification \\
POST & \code{/api/auth/email/verify/confirm/} & Non & Marque l'email comme vérifié \\
POST & \code{/api/auth/2fa/setup/} & Oui & Retourne le secret TOTP et l'URL \emph{otpauth} \\
POST & \code{/api/auth/2fa/confirm/} & Oui & Active la double authentification \\
POST & \code{/api/auth/2fa/disable/} & Oui & Désactive la double authentification \\
GET & \code{/api/auth/sessions/} & Oui & Liste les sessions actives \\
DELETE & \code{/api/auth/sessions/\{id\}/} & Oui & Révoque une session précise \\
\end{longtable}

\subsection{Extrait commenté — émission des jetons}

Chaque jeton de rafraîchissement est \emph{matérialisé} en base sous forme de
session. C'est ce qui permet à l'utilisateur de voir la liste de ses appareils
connectés et d'en révoquer un individuellement — ce qu'un JWT, sans état par
nature, ne permettrait pas seul.

\begin{lstlisting}[language=Python, caption={\code{auth\_app/services.py} — émission d'un couple de jetons et création de la session}]
def issue_auth_tokens(user, request):
    # SimpleJWT genere le couple acces/rafraichissement, signe en HS256
    # avec JWT_SIGNING_KEY -- le secret partage par TOUS les services.
    refresh = RefreshToken.for_user(user)

    expires_at = timezone.datetime.fromtimestamp(
        int(refresh["exp"]), tz=timezone.get_current_timezone()
    )

    # On MATERIALISE le jeton de rafraichissement en base. Le JWT reste sans
    # etat pour la validation (aucun appel reseau), mais on garde une trace
    # cote serveur pour pouvoir REVOQUER une session precise : c'est ce qui
    # alimente l'ecran "Sessions actives" et le bouton "Deconnecter partout".
    session = AuthSession.objects.create(
        user=user,
        refresh_jti=str(refresh["jti"]),   # identifiant unique du jeton
        ip_address=client_ip(request),
        user_agent=user_agent(request),
        expires_at=expires_at,
    )

    # On injecte l'id de session DANS le jeton : au rafraichissement, on peut
    # ainsi verifier que la session n'a pas ete revoquee entre-temps.
    refresh["session_id"] = session.id
    access = refresh.access_token

    return {
        "access": str(access),
        "refresh": str(refresh),
        "session_id": session.id,
    }
\end{lstlisting}

\subsection{Extrait commenté — publication d'événement}

\begin{lstlisting}[language=Python, caption={\code{auth\_app/events.py} — publication transactionnelle sur RabbitMQ}]
EXCHANGE_NAME = os.environ.get("RABBITMQ_EXCHANGE", "flowforge.events")
PRODUCER = "auth-service"

def publish_domain_event(event_type, resource_type, resource_id,
                         actor_user_id=None, data=None):
    # Enveloppe COMMUNE a tous les services, quel que soit le langage.
    # C'est ce contrat unique qui permet a un consommateur Java de lire
    # sans effort un evenement emis par un producteur Python.
    payload = {
        "eventId": str(uuid.uuid4()),
        "eventType": event_type,              # sert AUSSI de routing key
        "occurredAt": timezone.now().isoformat(),
        "producer": PRODUCER,
        "actorUserId": actor_user_id,
        "resource": {"type": resource_type, "id": resource_id},
        "data": data or {},
    }

    def _publish():
        rabbitmq_url = os.environ.get("RABBITMQ_URL")
        if not rabbitmq_url:
            return  # environnement sans bus (tests) : on n'echoue pas

        try:
            connection = pika.BlockingConnection(pika.URLParameters(rabbitmq_url))
            channel = connection.channel()
            # Idempotent : declarer un exchange deja existant ne fait rien.
            channel.exchange_declare(exchange=EXCHANGE_NAME,
                                     exchange_type="topic", durable=True)
            channel.basic_publish(
                exchange=EXCHANGE_NAME,
                routing_key=event_type,
                body=json.dumps(payload, default=str).encode("utf-8"),
                properties=pika.BasicProperties(
                    content_type="application/json",
                    delivery_mode=2,          # message PERSISTANT sur disque
                ),
            )
            connection.close()
        except Exception:
            # Publication "au mieux" : si le bus est indisponible, la requete
            # HTTP de l'utilisateur reussit quand meme. On degrade l'effet de
            # bord, jamais l'action principale.
            logger.exception("Failed to publish domain event %s", event_type)

    # >>> LE POINT ESSENTIEL <<<
    # On ne publie qu'APRES la validation de la transaction. Si l'INSERT est
    # annule, aucun evenement ne part : on ne signale jamais la creation d'un
    # utilisateur qui n'existe pas en base.
    transaction.on_commit(_publish)
\end{lstlisting}

% -----------------------------------------------------------------------------
\section{User-Service}
% -----------------------------------------------------------------------------

\subsection{Rôle}

Le \code{user-service} détient les \emph{métadonnées} de profil : nom affiché,
biographie, photo, fuseau horaire, langue, préférences. Il ne connaît ni les
mots de passe ni les sessions — cette séparation est volontaire : le service le
plus exposé fonctionnellement (celui qu'on modifie le plus souvent) est aussi
celui qui manipule le moins de données sensibles.

\subsection{Pile technique}

Django~5.2 avec Django REST Framework, comme l'auth-service — mais sur une base
\emph{séparée}. Le choix d'un second service Django plutôt qu'une simple
application supplémentaire dans l'auth-service est délibéré : profils et
authentification ont des cycles de vie, des besoins de mise à l'échelle et des
niveaux de sensibilité différents. La publication et la consommation d'événements
s'appuient sur la bibliothèque \code{pika}.

\subsection{Structure du projet}

\begin{lstlisting}[language=bash,numbers=none]
server/user-service/
  users_django/           # configuration du projet Django
    settings.py           #   BD, stockage Azure, RabbitMQ
    cors.py               #   middleware CORS ecrit a la main
  profiles/
    models.py             #   Profile
    serializers.py        #   validation des entrees (DRF)
    services.py           #   logique metier (photo, preferences)
    views.py              #   points d'entree HTTP
    storage.py            #   depot des photos vers Azure Blob
    events.py             #   publication RabbitMQ (user.profile.*)
    management/commands/
      consume_events.py   #   CONSOMMATEUR de user.registered
    migrations/
  entrypoint.sh           # lance le serveur HTTP + le consommateur
  Dockerfile
\end{lstlisting}

\subsection{Modèle de données}

Une seule table, \code{profiles\_profile} (voir aussi la
figure~\ref{fig:modele-auth}). Le champ pivot est \code{auth\_user\_id} : c'est
la copie de l'identifiant émis par l'auth-service, sans contrainte de clé
étrangère puisqu'il pointe vers une \emph{autre} base. Les autres colonnes sont
purement descriptives : \code{username}, \code{email}, \code{biography},
\code{timezone}, \code{language}, \code{preferences} (JSON libre) et les
métadonnées de la photo (\code{profile\_picture\_key} désigne l'objet stocké
dans Azure Blob, pas les octets eux-mêmes).

\subsection{Particularité : provisionnement automatique par événement}

Le profil n'est pas créé par un appel HTTP depuis l'auth-service : il est créé
par la \emph{consommation} de l'événement \code{user.registered}. Le
consommateur tourne comme un processus d'arrière-plan à côté du serveur HTTP,
lancé par le point d'entrée du conteneur.

\begin{lstlisting}[language=bash, caption={\code{entrypoint.sh} — le serveur HTTP et le consommateur d'événements dans le même conteneur}]
#!/bin/sh
set -e
python manage.py migrate --noinput

# Le consommateur RabbitMQ tourne en TACHE DE FOND...
python manage.py consume_events &

# ...pendant que le serveur HTTP occupe le processus principal.
exec python manage.py runserver 0.0.0.0:8000
\end{lstlisting}

\begin{lstlisting}[language=Python, caption={\code{profiles/management/commands/consume\_events.py} — déclaration de la topologie et du rebut}]
# Le consommateur declare LUI-MEME toute la topologie dont il depend.
# Il peut donc demarrer avant, apres, ou sans le producteur : l'ordre de
# demarrage des conteneurs n'a aucune importance.
channel.exchange_declare(exchange=exchange, exchange_type="topic", durable=True)

# Exchange de rebut (dead-letter) : y atterrissent les messages definitivement
# en echec, au lieu d'etre perdus ou rejoues indefiniment.
channel.exchange_declare(exchange=dlx, exchange_type="direct", durable=True)

channel.queue_declare(queue=queue, durable=True, arguments={
    "x-dead-letter-exchange": dlx,
    "x-dead-letter-routing-key": dlq_routing_key,
})
channel.queue_declare(queue=dlq, durable=True)

# On ne s'abonne QU'A la cle qui nous concerne. Les centaines d'autres
# evenements qui transitent par l'exchange ne nous sont jamais livres.
channel.queue_bind(queue=queue, exchange=exchange, routing_key="user.registered")

# Un seul message non acquitte a la fois : on ne prefetch pas un lot que
# l'on serait incapable de traiter en cas de redemarrage.
channel.basic_qos(prefetch_count=1)
\end{lstlisting}

\subsection{Points d'entrée exposés}

\begin{table}[H]
\centering
\small
\begin{tabularx}{\linewidth}{@{}p{1.5cm}p{5.4cm}p{1.2cm}X@{}}
\toprule
\textbf{Méth.} & \textbf{Chemin} & \textbf{Auth} & \textbf{Rôle} \\
\midrule
GET & \code{/api/profiles/} & Oui & Profil de l'utilisateur courant \\
PATCH & \code{/api/profiles/\{id\}/} & Oui & Mise à jour du profil \\
GET & \code{/api/profiles/public/\{username\}/} & Non & Profil public \\
POST & \code{/api/profiles/picture/} & Oui & Téléversement de la photo \\
PUT & \code{/api/profiles/picture/} & Oui & Remplacement de la photo \\
DELETE & \code{/api/profiles/picture/} & Oui & Suppression de la photo \\
\bottomrule
\end{tabularx}
\caption{Points d'entrée du User-Service}
\end{table}

La photo de profil est validée côté serveur (types \code{image/jpeg},
\code{image/png}, \code{image/webp} ; 2~Mo maximum) puis déposée dans un
stockage objet Azure Blob. La base ne contient que la \emph{clé} de l'objet et
ses métadonnées : elle ne grossit pas avec le volume des images.

% -----------------------------------------------------------------------------
\section{TODO-Service (task-service)}
% -----------------------------------------------------------------------------

\subsection{Rôle}

C'est le cœur métier de la plateforme. Il possède les tâches et tout ce qui
gravite autour : sous-tâches, étiquettes, assignations, dépendances, récurrence,
rappels, colonnes kanban, filtres enregistrés et journal d'activité.

\subsection{Pile technique}

NestJS~11 avec TypeORM. NestJS apporte une structure modulaire explicite
(modules, contrôleurs, services, \emph{guards}, \emph{pipes}) et une validation
déclarative des entrées via \code{class-validator}, activée globalement :

\begin{lstlisting}[language=TypeScript, caption={\code{src/main.ts} — validation stricte des entrées}]
app.useGlobalPipes(new ValidationPipe({
  whitelist: true,            // supprime les proprietes non declarees dans le DTO
  forbidNonWhitelisted: true, // ...et REJETTE la requete (400) si elle en contient
  transform: true,            // convertit les types (ex. "3" -> 3)
}));
\end{lstlisting}

\code{forbidNonWhitelisted} est un choix de sécurité : un client qui tenterait
d'injecter un champ non prévu (\code{"owner\_user\_id": 1}) reçoit une erreur
400 au lieu de voir son champ silencieusement ignoré.

\subsection{Structure du projet}

C'est le service le plus riche : son domaine est découpé en modules NestJS
distincts (tâches, statuts, colonnes kanban, filtres enregistrés), chacun avec
son contrôleur et son service.

\begin{lstlisting}[language=bash,numbers=none]
server/task-service/
  src/
    main.ts                 # amorcage, ValidationPipe globale, CORS
    app.module.ts           # assemblage des modules + TypeORM
    auth/
      jwt-auth.guard.ts      #   validation locale du JWT (HS256)
      current-user.decorator.ts
    tasks/
      tasks.controller.ts    #   ~30 routes (taches + collaboration)
      tasks.service.ts       #   logique metier
      deadline-scanner.service.ts  #   cron des echeances (bonus)
      dto/task.dto.ts        #   DTO valides par class-validator
      entities/              #   Task, Subtask, Tag, Assignment...
    events/
      domain-event.publisher.ts    #   publication RabbitMQ (task.*)
    migrations/              # migrations TypeORM
  Dockerfile
\end{lstlisting}

\subsection{Modèle de données}

\mermaidfig{09-modele-task}{Modèle de données du TODO-Service (11 entités)}{modele-task}{1.0}

\begin{choix}[Le statut : d'une table de jointure à un type énuméré natif]
Dans une première version, le statut d'une tâche était une clé étrangère vers
une table \code{task\_statuses} configurable par utilisateur. Nous l'avons
\textbf{dénormalisé} en un type énuméré natif PostgreSQL
(\code{backlog}, \code{todo}, \code{in\_progress}, \code{review}, \code{done},
\code{archived}). Motif : la flexibilité offerte n'était jamais utilisée, alors
qu'elle imposait une jointure sur \emph{chaque} lecture de tâche et rendait
impossible tout filtrage efficace par statut. Les colonnes
\code{completed\_at} et \code{archived\_at} sont désormais \emph{dérivées} du
statut, ce qui supprime une classe entière d'incohérences (une tâche marquée
« terminée » sans date de complétion, par exemple).
\end{choix}

\begin{lstlisting}[language=SQL, caption={\code{src/migrations/…-DirectTaskStatusAndActivity.ts} — migration de dénormalisation}]
-- 1. Creation du type enumere natif
CREATE TYPE tasks_status_enum AS ENUM
  ('backlog','todo','in_progress','review','done','archived');

-- 2. Ajout de la colonne, avec une valeur par defaut sure
ALTER TABLE tasks ADD COLUMN IF NOT EXISTS
  status tasks_status_enum NOT NULL DEFAULT 'todo';

-- 3. Reprise des donnees existantes depuis l'ancienne table (non montre),
--    PUIS seulement suppression de l'ancien schema.
ALTER TABLE tasks DROP COLUMN IF EXISTS status_id;
DROP TABLE IF EXISTS task_statuses;

-- 4. Nouveau journal d'activite : trace horodatee de chaque changement
CREATE TABLE IF NOT EXISTS task_activities (
  id            SERIAL PRIMARY KEY,
  task_id       INT NOT NULL,
  actor_user_id INT NOT NULL,
  type          VARCHAR(80) NOT NULL,   -- status_changed, assignment_added...
  from_value    VARCHAR(200),           -- valeur AVANT
  to_value      VARCHAR(200),           -- valeur APRES
  message       VARCHAR(500) NOT NULL,
  created_at    TIMESTAMPTZ NOT NULL DEFAULT NOW(),
  CONSTRAINT fk_task_activities_task
    FOREIGN KEY (task_id) REFERENCES tasks(id) ON DELETE CASCADE
);
CREATE INDEX IF NOT EXISTS idx_task_activities_task_id ON task_activities(task_id);
\end{lstlisting}

\subsection{Points d'entrée exposés}

\begin{longtable}{@{}p{1.4cm}p{7.2cm}p{5.9cm}@{}}
\caption{Principaux points d'entrée du TODO-Service}\label{tab:ep-task}\\
\toprule
\textbf{Méth.} & \textbf{Chemin} & \textbf{Rôle} \\
\midrule
\endfirsthead
\toprule
\textbf{Méth.} & \textbf{Chemin} & \textbf{Rôle} \\
\midrule
\endhead
\bottomrule
\endlastfoot
GET & \code{/api/tasks/} & Liste filtrée (catégorie, statut, priorité…) \\
POST & \code{/api/tasks/} & Création d'une tâche \\
GET & \code{/api/tasks/\{id\}} & Détail d'une tâche \\
PATCH & \code{/api/tasks/\{id\}} & Mise à jour partielle \\
DELETE & \code{/api/tasks/\{id\}} & Suppression logique \\
POST & \code{/api/tasks/\{id\}/complete} & Marque la tâche terminée \\
POST & \code{/api/tasks/\{id\}/archive} & Archive la tâche \\
GET & \code{/api/tasks/\{id\}/permissions} & Droits de l'appelant sur la tâche \\
GET & \code{/api/tasks/\{id\}/activity} & Journal d'activité \\
POST & \code{/api/tasks/\{id\}/assignments} & Assigne un utilisateur \\
DELETE & \code{/api/tasks/\{id\}/assignments/\{userId\}} & Retire une assignation \\
POST & \code{/api/tasks/\{id\}/subtasks} & Crée une sous-tâche \\
PATCH & \code{/api/tasks/\{id\}/subtasks/\{sid\}} & Met à jour une sous-tâche \\
DELETE & \code{/api/tasks/\{id\}/subtasks/\{sid\}} & Supprime une sous-tâche \\
POST & \code{/api/tasks/\{id\}/tags} & Ajoute une étiquette \\
DELETE & \code{/api/tasks/\{id\}/tags/\{tagId\}} & Retire une étiquette \\
POST & \code{/api/tasks/\{id\}/dependencies} & Ajoute une dépendance \\
PUT & \code{/api/tasks/\{id\}/recurrence} & Définit la récurrence \\
POST & \code{/api/tasks/\{id\}/reminders} & Crée un rappel \\
GET/POST & \code{/api/tasks/kanban-columns} & Colonnes kanban \\
GET/POST & \code{/api/tasks/saved-filters} & Filtres enregistrés \\
\end{longtable}

\subsection{Extrait commenté — le \emph{guard} JWT}

Ce \emph{guard} est le même dans les quatre services NestJS. Il illustre
concrètement l'authentification sans état : aucune requête réseau, aucun accès
base, une simple vérification de signature.

\begin{lstlisting}[language=TypeScript, caption={\code{src/auth/jwt-auth.guard.ts} — validation locale du jeton}]
@Injectable()
export class JwtAuthGuard implements CanActivate {
  constructor(
    private readonly jwt: JwtService,
    private readonly config: ConfigService,
  ) {}

  async canActivate(context: ExecutionContext): Promise<boolean> {
    const request = context.switchToHttp().getRequest<Request>();
    const header = request.headers.authorization;

    if (!header?.startsWith('Bearer ')) {
      throw new UnauthorizedException('Missing bearer token');
    }
    const token = header.slice('Bearer '.length).trim();

    let payload: AccessTokenPayload;
    try {
      // Verification PUREMENT LOCALE : on recalcule la signature HS256 avec le
      // secret partage. Aucun appel a l'auth-service, donc aucune latence
      // reseau et aucun point de defaillance unique sur ce chemin.
      payload = await this.jwt.verifyAsync<AccessTokenPayload>(token, {
        secret: this.config.getOrThrow<string>('JWT_SIGNING_KEY'),
        algorithms: ['HS256'],   // on FIXE l'algorithme : se protege de
                                 // l'attaque "alg: none" / substitution
      });
    } catch {
      throw new UnauthorizedException('Invalid or expired token');
    }

    const userId = Number(payload.user_id);
    if (!Number.isFinite(userId)) {
      throw new UnauthorizedException('Token has no user_id claim');
    }

    // On attache l'utilisateur A LA REQUETE. On conserve aussi le jeton BRUT :
    // certains services (statistics, comment) doivent le REJOUER lors de leurs
    // appels sortants, afin d'heriter exactement des droits de l'appelant.
    (request as Request & { user: AuthUser }).user = { userId, token };
    return true;
  }
}
\end{lstlisting}

% -----------------------------------------------------------------------------
\section{Category-Service}
% -----------------------------------------------------------------------------

\subsection{Rôle}

Il possède les catégories (Travail, Personnel, Urgent…), leurs membres, leurs
rôles et leurs invitations. Une tâche appartient à une catégorie ; une catégorie
appartient à un espace de travail.

\subsection{Pile technique}

NestJS~11 avec TypeORM, comme le task-service et le comment-service. Les trois
services NestJS partagent volontairement la même ossature (module \code{auth}
avec le \emph{guard} JWT, module \code{events} avec le publieur RabbitMQ,
\code{ValidationPipe} globale), ce qui rend le code prévisible d'un service à
l'autre. La particularité du category-service est qu'il ne prend \emph{aucune}
décision d'autorisation seul : il appelle le workspace-service en HTTP pour
connaître les droits hérités (voir l'extrait plus bas).

\subsection{Structure du projet}

\begin{lstlisting}[language=bash,numbers=none]
server/category-service/
  src/
    main.ts                 # amorcage, ValidationPipe globale, CORS
    app.module.ts           # assemblage + TypeORM (base categories)
    auth/
      jwt-auth.guard.ts      #   validation locale du JWT (HS256)
      current-user.decorator.ts
    categories/
      categories.controller.ts   #   CRUD + membres + invitations
      categories.service.ts      #   logique metier + droits herites
      dto/category.dto.ts
      entities/category.entity.ts  #   Category, CategoryMember,
                                   #   CategoryInvitation (3 entites)
    events/
      domain-event.publisher.ts    #   publication RabbitMQ (category.*)
  Dockerfile
\end{lstlisting}

\subsection{Modèle de données}

\mermaidfig{10-modele-category-comment}{Modèle de données des services Category et Comment}{modele-cat}{0.95}

Trois tables : \code{categories}, \code{category\_members} (adhésions explicites
avec un rôle) et \code{category\_invitations} (invitations par jeton). Comme
partout, il n'y a aucune clé étrangère vers les autres services :
\code{workspace\_id}, \code{owner\_user\_id} et \code{author\_user\_id} sont de
simples entiers. On notera le champ \code{workspace\_name\_snapshot} :
c'est une \emph{copie} du nom de l'espace de travail au moment de la création,
qui évite un appel réseau à chaque affichage — une dénormalisation assumée.

\subsection{Points d'entrée exposés}

\begin{longtable}{@{}p{1.4cm}p{7.6cm}p{5.5cm}@{}}
\caption{Points d'entrée du Category-Service}\label{tab:ep-cat}\\
\toprule
\textbf{Méth.} & \textbf{Chemin} & \textbf{Rôle} \\
\midrule
\endfirsthead
\toprule
\textbf{Méth.} & \textbf{Chemin} & \textbf{Rôle} \\
\midrule
\endhead
\bottomrule
\endlastfoot
GET & \code{/api/categories/} & Liste des catégories visibles \\
POST & \code{/api/categories/} & Création \\
GET & \code{/api/categories/\{id\}} & Détail \\
PATCH & \code{/api/categories/\{id\}} & Mise à jour (admin) \\
DELETE & \code{/api/categories/\{id\}} & Suppression logique \\
POST & \code{/api/categories/\{id\}/archive} & Archivage \\
GET & \code{/api/categories/\{id\}/permissions/me} & Droits de l'appelant \\
GET & \code{/api/categories/\{id\}/members} & Liste des membres \\
PATCH & \code{/api/categories/\{id\}/members/\{uid\}} & Changement de rôle \\
DELETE & \code{/api/categories/\{id\}/members/\{uid\}} & Retrait d'un membre \\
GET & \code{/api/categories/\{id\}/invitations} & Liste des invitations \\
POST & \code{/api/categories/\{id\}/invitations} & Invitation par email \\
POST & \code{/api/categories/invitations/\{token\}/accept} & Acceptation \\
POST & \code{/api/categories/invitations/\{token\}/reject} & Refus \\
\end{longtable}

\subsection{Extrait commenté — héritage des rôles}

C'est la règle métier la plus subtile du service : les droits sur une catégorie
sont \emph{hérités} de l'espace de travail, sauf si une adhésion explicite les
surcharge.

\begin{lstlisting}[language=TypeScript, caption={\code{src/categories/categories.service.ts} — résolution des droits}]
private async resolvePermission(user, category): Promise<CategoryPermission> {
  // Deux sources de droits, interrogees EN PARALLELE :
  //  1. les droits herites de l'espace de travail (appel HTTP au workspace-service)
  //  2. une eventuelle adhesion EXPLICITE a cette categorie (base locale)
  const [workspacePermission, member] = await Promise.all([
    this.workspacePermissionOrNull(user, category.workspace_id),
    this.membersRepo.findOne({
      where: {
        category_id: category.id,
        user_id: user.userId,
        removed_at: IsNull(),   // un membre "retire" ne compte plus
      },
    }),
  ]);

  const workspaceRole = workspacePermission?.role ?? null;

  // L'adhesion EXPLICITE prime sur le role herite. C'est ce qui permet
  // d'inviter quelqu'un sur UNE SEULE categorie sans lui donner acces a
  // tout l'espace de travail.
  const categoryRole = member?.role ?? this.derivedCategoryRole(workspaceRole);

  const isAdmin =
    workspacePermission?.canAdmin === true ||
    categoryRole === CategoryRole.CATEGORY_OWNER ||
    categoryRole === CategoryRole.CATEGORY_ADMIN;

  const canContribute = isAdmin || categoryRole === CategoryRole.CONTRIBUTOR;
  const canView = workspacePermission?.canView === true || Boolean(member);
  // ...
}

// Table de correspondance des roles herites.
private derivedCategoryRole(role: WorkspacePermission['role']): CategoryRole | null {
  if (role === 'OWNER')  return CategoryRole.CATEGORY_OWNER;
  if (role === 'ADMIN')  return CategoryRole.CATEGORY_ADMIN;
  if (role === 'MEMBER') return CategoryRole.CONTRIBUTOR;
  if (role === 'VIEWER') return CategoryRole.VIEWER;
  return null;   // non-membre : aucun droit herite
}
\end{lstlisting}

% -----------------------------------------------------------------------------
\section{Comment-Service}
% -----------------------------------------------------------------------------

\subsection{Rôle}

Il possède les commentaires, et \emph{uniquement} les commentaires. C'est le
service qui illustre le mieux la discipline du découpage : il ne stocke aucune
donnée de tâche, aucune donnée d'utilisateur, et ne réimplémente aucun contrôle
d'accès.

\begin{choix}[Ne jamais dupliquer une règle de sécurité]
Pour savoir si un utilisateur a le droit de lire ou d'écrire un commentaire, la
question à poser est en réalité : « a-t-il le droit de voir la \emph{tâche}
commentée ? ». Cette règle appartient au \code{task-service}. Le
\code{comment-service} ne la réécrit donc pas : il l'\textbf{interroge}, en
appelant \code{GET /api/tasks/\{id\}/permissions} avec le jeton de l'appelant.
Le coût est un appel HTTP supplémentaire ; le bénéfice est qu'il n'existe qu'un
seul endroit au monde où la règle de visibilité d'une tâche est écrite. Elle ne
peut donc pas diverger.
\end{choix}

\subsection{Pile technique}

NestJS~11 avec TypeORM, sur sa propre base \code{comments}. Ce service a été
\emph{extrait} du task-service : les commentaires y étaient auparavant un module
interne. Les isoler dans un service dédié, avec sa propre base, permet de les
faire évoluer (et monter en charge) indépendamment des tâches.

\subsection{Structure du projet}

C'est le plus petit des services NestJS — la meilleure preuve qu'un microservice
peut rester minuscule quand son périmètre est étroit.

\begin{lstlisting}[language=bash,numbers=none]
server/comment-service/
  src/
    main.ts                 # amorcage, ValidationPipe globale, CORS
    app.module.ts           # assemblage + TypeORM (base comments)
    auth/
      jwt-auth.guard.ts      #   validation locale du JWT (HS256)
      current-user.decorator.ts
    comments/
      comments.controller.ts     #   5 routes CRUD
      comments.service.ts        #   appelle task-service pour les droits
      dto/comment.dto.ts
      entities/comment.entity.ts #   une seule entite
    events/
      domain-event.publisher.ts  #   publication RabbitMQ (task.comment.*)
  Dockerfile
\end{lstlisting}

\subsection{Modèle de données}

Une seule table (\code{comments}), sans aucune clé étrangère : \code{task\_id},
\code{category\_id}, \code{workspace\_id} et \code{author\_user\_id} désignent
des lignes situées dans d'\emph{autres} bases.

\begin{lstlisting}[language=TypeScript, caption={\code{src/comments/entities/comment.entity.ts}}]
@Entity({ name: 'comments' })
export class Comment {
  @PrimaryGeneratedColumn()
  id: number;

  // Reference vers une tache detenue par le task-service, dans une AUTRE base.
  // C'est un simple entier indexe : aucune contrainte FK n'est possible, et
  // c'est le prix assume du decouplage. L'integrite est maintenue par les
  // evenements (task.deleted -> suppression des commentaires associes).
  @Column({ type: 'int' })
  @Index()
  task_id: number;

  @Column({ type: 'int', nullable: true })
  @Index()
  category_id: number | null;

  @Column({ type: 'int', nullable: true })
  @Index()
  workspace_id: number | null;

  @Column({ type: 'int' })
  @Index()
  author_user_id: number;

  @Column({ type: 'text' })
  body: string;

  // Suppression LOGIQUE : un commentaire supprime reste en base (tracabilite)
  // mais disparait de toutes les lectures et de l'index de recherche.
  @Column({ type: 'timestamptz', nullable: true })
  deleted_at: Date | null;

  @Column({ type: 'int', nullable: true })
  deleted_by_user_id: number | null;

  @CreateDateColumn({ name: 'created_at' })
  created_at: Date;

  @UpdateDateColumn({ name: 'updated_at' })
  updated_at: Date;
}
\end{lstlisting}

\subsection{Points d'entrée exposés}

\begin{table}[H]
\centering
\small
\begin{tabularx}{\linewidth}{@{}p{1.5cm}p{5.6cm}X@{}}
\toprule
\textbf{Méth.} & \textbf{Chemin} & \textbf{Rôle} \\
\midrule
GET & \code{/api/comments/?task\_id=\{id\}} & Commentaires d'une tâche \\
POST & \code{/api/comments/} & Ajout d'un commentaire \\
GET & \code{/api/comments/\{id\}} & Détail \\
PATCH & \code{/api/comments/\{id\}} & Modification (auteur uniquement) \\
DELETE & \code{/api/comments/\{id\}} & Suppression logique \\
\bottomrule
\end{tabularx}
\caption{Points d'entrée du Comment-Service}
\end{table}

\subsection{Extrait commenté — publication enrichie}

Le commentaire publié sur le bus embarque la liste des utilisateurs autorisés à
le voir (\code{viewerUserIds}). Cette liste sera recopiée telle quelle dans le
document Elasticsearch, ce qui permettra de filtrer les résultats de recherche
\emph{à l'intérieur du moteur} (voir section~\ref{sec:search}).

\begin{lstlisting}[language=TypeScript, caption={\code{src/comments/comments.service.ts} — enrichissement de l'événement}]
private async publish(eventType: string, comment: Comment,
                      actorUserId: number, task: TaskSnapshot) {
  await this.events.publish({
    eventType,                                    // task.comment.created/updated/deleted
    actorUserId,
    resource: { type: 'comment', id: comment.id },
    data: {
      ...this.searchableData(comment),
      title: task.title,

      // La LISTE DES AYANTS DROIT est calculee ICI, au moment de la
      // publication, et voyage AVEC l'evenement. Le search-service n'a donc
      // aucun appel a faire pour savoir qui peut voir ce commentaire : il
      // recopie cette liste dans le document indexe.
      viewerUserIds: [
        task.author_user_id,
        comment.author_user_id,
        ...(task.assignments ?? []).map((a) => a.assignee_user_id),
      ],
    },
  });
}
\end{lstlisting}

% -----------------------------------------------------------------------------
\section{Statistics-Service}
% -----------------------------------------------------------------------------

\subsection{Rôle}

Il calcule les métriques de productivité : taux de complétion, totaux,
répartitions par statut, catégorie et priorité, et tendance de complétion.

\subsection{Pile technique}

NestJS~11, mais \emph{sans} TypeORM ni base : c'est le seul service NestJS
réduit à un contrôleur, un service métier et le \emph{guard} JWT. Il ne conserve
que le module \code{auth} commun aux autres services NestJS ; tout le reste
(entités, dépôts, publieur d'événements) est absent, parce qu'inutile.

\subsection{Structure du projet}

\begin{lstlisting}[language=bash,numbers=none]
server/statistics-service/
  src/
    main.ts                 # amorcage, CORS
    app.module.ts           # AUCUN TypeOrmModule -- pas de base
    auth/
      jwt-auth.guard.ts      #   validation locale du JWT (HS256)
      current-user.decorator.ts
    statistics/
      statistics.controller.ts   #   une seule route : /overview
      statistics.service.ts      #   appel task-service + agregation
  Dockerfile
\end{lstlisting}

\subsection{Modèle de données}

\textbf{Aucun.} Ce service ne possède ni table, ni base de données — c'est
assumé et central à sa conception. L'encadré ci-dessous explique pourquoi.

\begin{choix}[Le seul service sans base de données — et c'est délibéré]
Un service de statistiques n'a besoin d'aucune donnée \emph{qui lui soit
propre} : il ne fait que \emph{regarder} celles des autres. Lui donner une base
signifierait dupliquer les tâches, puis assumer la désynchronisation
inévitable entre la copie et l'original.

Nous avons donc choisi d'en faire un \textbf{agrégateur sans état}. Il
\emph{rejoue} le jeton JWT de l'utilisateur auprès du \code{task-service}, et
n'obtient donc, par construction, que les tâches que cet utilisateur a le droit
de voir. \textbf{Aucun contrôle d'accès n'est réimplémenté} : il est
\emph{hérité}. L'agrégation se fait ensuite en mémoire.

C'est aussi la limite du choix : chaque appel refait une requête HTTP complète
vers le \code{task-service}. À très grande échelle, il faudrait basculer sur un
modèle de lecture CQRS alimenté par les événements \code{task.*} — toute
l'infrastructure d'événements nécessaire existe déjà.
\end{choix}

\mermaidfig{13-flux-statistiques}{Flux d'agrégation des statistiques}{stats}{0.95}

\subsection{Point d'entrée exposé}

\begin{table}[H]
\centering
\small
\begin{tabularx}{\linewidth}{@{}p{1.3cm}p{5.2cm}X@{}}
\toprule
\textbf{Méth.} & \textbf{Chemin} & \textbf{Filtres} \\
\midrule
GET & \code{/api/statistics/overview} & \code{workspace\_id}, \code{category\_id}, \code{priority}, \code{status} \\
\bottomrule
\end{tabularx}
\caption{Point d'entrée du Statistics-Service}
\end{table}

Réponse : \code{totalTasks}, \code{completedTasks}, \code{completionRate},
\code{byStatus}, \code{byCategory}, \code{byPriority}, \code{trends}.

\subsection{Extrait commenté — agrégation}

\begin{lstlisting}[language=TypeScript, caption={\code{src/statistics/statistics.service.ts} — délégation des droits puis agrégation}]
async overview(user: { token: string }, filters: OverviewFilters) {
  const tasks = (await this.fetchVisibleTasks(user, filters))
    .filter((task) => this.matches(task, filters));

  const total = tasks.length;
  const completed = tasks.filter(
    (task) => task.status === 'done' || Boolean(task.completed_at),
  ).length;

  return {
    totalTasks: total,
    completedTasks: completed,
    // Taux arrondi a une decimale ; 0 si aucune tache (evite la division par 0)
    completionRate: total ? Math.round((completed / total) * 1000) / 10 : 0,
    byStatus:   this.distribution(tasks, (t) => t.status ?? 'unknown'),
    byCategory: this.distribution(tasks, (t) => String(t.category_id ?? 'uncategorized')),
    byPriority: this.distribution(tasks, (t) => t.priority ?? 'normal'),
    trends:     this.dailyCompletedTrend(tasks),
  };
}

private async fetchVisibleTasks(user: { token: string }, filters: OverviewFilters) {
  const params = new URLSearchParams();
  if (filters.workspaceId) params.set('workspace_id', String(filters.workspaceId));
  if (filters.categoryId)  params.set('category_id',  String(filters.categoryId));
  // ...

  const response = await fetch(`${this.taskBaseUrl()}/api/tasks/?${params}`, {
    // >>> LE POINT CLE <<<
    // On REJOUE le jeton de l'utilisateur. Le task-service applique donc SES
    // regles de visibilite habituelles, et ne nous renvoie que les taches que
    // CET utilisateur peut voir. Les statistiques ne peuvent structurellement
    // pas fuiter la tache d'un autre : le controle d'acces est herite, pas
    // reimplemente.
    headers: { Authorization: `Bearer ${user.token}` },
  });

  if (response.status === 401 || response.status === 403) {
    throw new ForbiddenException('Task statistics access denied.');
  }
  return (await response.json()) as TaskSnapshot[];
}
\end{lstlisting}

% -----------------------------------------------------------------------------
\section{Notification-Service \textnormal{\itshape(bonus)}}
\label{sec:notif}
% -----------------------------------------------------------------------------

\subsection{Rôle}

Le service optionnel valorisé en bonus. Il détecte les tâches dont l'échéance
est dépassée ou proche (moins de 24~h), génère des alertes par utilisateur, et
les pousse en temps réel dans l'interface.

\subsection{Pile technique}

Spring Boot~3.4 / Java~21, \textbf{MySQL~8}, WebSocket natif Spring, consommation
RabbitMQ via Spring AMQP. Le choix de Spring pour ce service tient à ses briques
intégrées : \code{spring-boot-starter-websocket} pour le temps réel,
Spring AMQP pour la consommation d'événements avec réessais, et Flyway pour le
versionnement du schéma MySQL.

\subsection{Structure du projet}

Le code est organisé par responsabilité technique (\emph{package by layer}),
convention idiomatique de Spring Boot.

\begin{lstlisting}[language=bash,numbers=none]
server/notification-service/src/main/
  java/com/flowforge/notification/
    config/
      RabbitEventConfig.java     # file durable + DLQ + reessais
      SecurityConfig.java        # validation JWT
    controller/
      NotificationController.java  # 5 routes REST
    event/
      RabbitNotificationEventListener.java  # @RabbitListener
      DomainNotificationEventMapper.java    # 1 evenement -> N notifs
    model/
      Notification.java          # entite JPA
      NotificationType.java      # enumeration des types
    repository/
      NotificationRepository.java
    service/
      NotificationService.java   # creation idempotente
    websocket/
      NotificationWebSocketHandler.java    # socket authentifiee par JWT
      NotificationWebSocketPublisher.java  # push temps reel
  resources/
    application.yml
    db/migration/                # scripts Flyway (schema MySQL)
\end{lstlisting}

\subsection{Modèle de données}

Une seule table, \code{notifications} (détaillée sur la
figure~\ref{fig:modele-ws}), sur MySQL~8. La colonne décisive est la contrainte
d'unicité \code{(source\_event\_id, recipient\_user\_id, type)} : c'est elle qui
rend la création idempotente et absorbe les redistributions de RabbitMQ (voir
l'extrait plus bas). Les champs \code{related\_resource\_id} et
\code{related\_resource\_type} forment un pointeur polymorphe vers la ressource
d'origine (une tâche, un espace de travail…), et \code{deleted\_at} matérialise
la suppression logique.

\subsection{Détection des échéances}

Une tâche planifiée (\emph{cron}) balaie chaque heure les tâches dont
l'échéance tombe dans les 24 prochaines heures et qui ne sont pas terminées,
puis publie un événement \code{task.deadline\_approaching} par tâche
concernée.

\mermaidfig{07-sequence-deadline}{Détection des échéances proches et notification temps réel}{deadline}{0.95}

\begin{lstlisting}[language=TypeScript, caption={\code{task-service} — l'ordonnanceur qui déclenche la chaîne de notification}]
@Injectable()
export class DeadlineScannerService {
  // Balayage horaire. L'expression cron est configurable (DEADLINE_SCAN_CRON).
  @Cron(process.env.DEADLINE_SCAN_CRON ?? CronExpression.EVERY_HOUR)
  async scanApproachingDeadlines(): Promise<void> {
    const now = new Date();
    const horizon = new Date(now.getTime() + this.windowHours * 3600 * 1000);

    const dueSoon = await this.tasksRepo.find({
      where: {
        due_at: Between(now, horizon),  // echeance dans la fenetre
        completed_at: IsNull(),         // pas deja terminee
        deleted_at: IsNull(),           // pas supprimee
      },
      relations: ['assignments'],
    });

    for (const task of dueSoon) {
      // On publie UN evenement par tache. Le task-service ne sait pas -- et
      // n'a pas besoin de savoir -- ce qui en sera fait : notification in-app,
      // email, webhook... Ajouter un nouveau consommateur ne changera pas
      // une ligne de ce fichier.
      await this.events.publish({
        eventType: 'task.deadline_approaching',
        actorUserId: null,                       // declenche par le systeme
        resource: { type: 'task', id: task.id },
        data: {
          taskId: task.id,
          title: task.title,
          dueAt: task.due_at,
          workspaceId: task.workspace_id,
          // Destinataires : les assignes de la tache
          viewerUserIds: task.assignments.map((a) => a.assignee_user_id),
        },
      });
    }
  }
}
\end{lstlisting}

\subsection{Extrait commenté — idempotence et éventail de destinataires}

Un même événement peut concerner plusieurs personnes. Le service le transforme
en \emph{N} notifications — une par destinataire — en prenant soin d'exclure
l'auteur de l'action.

\begin{lstlisting}[language=Java, caption={\code{DomainNotificationEventMapper.java} — un événement, N notifications}]
private List<Long> recipients(Map<String, Object> data, Long actorUserId,
                              String... keys) {
    // LinkedHashSet : deduplique tout en preservant l'ordre.
    Set<Long> recipients = new LinkedHashSet<>();
    for (var key : keys) {
        recipients.addAll(numbers(data, key));
    }

    // >>> On RETIRE l'acteur de la liste. <<<
    // Sans cela, celui qui s'assigne une tache a lui-meme recevrait une
    // notification lui annoncant qu'il vient de s'assigner une tache.
    // Le garde-fou "size() > 1" evite de vider la liste dans le cas ou
    // l'acteur est le SEUL destinataire legitime.
    if (actorUserId != null && recipients.size() > 1) {
        recipients.remove(actorUserId);
    }
    return List.copyOf(recipients);
}
\end{lstlisting}

\begin{lstlisting}[language=Java, caption={\code{NotificationService.java} — idempotence face aux redistributions RabbitMQ}]
// RabbitMQ garantit "au moins une fois", PAS "exactement une fois" : un
// message PEUT etre redistribue (redemarrage du consommateur, acquittement
// perdu...). Sans protection, l'utilisateur verrait la meme notification en
// double.
//
// On s'appuie sur une contrainte d'unicite (source_event_id, recipient, type).
if (sourceEventId != null && !sourceEventId.isBlank()) {
    var existing = notifications.findBySourceEventIdAndRecipientUserIdAndType(
        sourceEventId, recipientUserId, type);
    if (existing.isPresent()) {
        return new NotificationCreation(toResponse(existing.get()), false);
        //                                    "created = false" ^^^^^
        // => le WebSocket ne repoussera PAS la notification. Le rejeu est
        //    totalement transparent pour l'utilisateur final.
    }
}

try {
    // ... insertion
} catch (DataIntegrityViolationException exc) {
    // Deux messages traites EN MEME TEMPS peuvent passer le test ci-dessus
    // tous les deux. La contrainte d'unicite en base tranche : on rattrape
    // l'exception et on renvoie la ligne gagnante.
    return existing(...);
}
\end{lstlisting}

\subsection{Temps réel : WebSocket authentifié par le jeton}

\begin{lstlisting}[language=Java, caption={\code{NotificationWebSocketHandler.java} — authentification d'une socket}]
// PROBLEME : un navigateur ne peut PAS ajouter d'en-tete "Authorization" a une
// connexion WebSocket (l'API WebSocket du navigateur ne le permet pas).
// SOLUTION : le jeton voyage en parametre de requete, et la connexion est
// refusee immediatement s'il est absent ou invalide.
var token = UriComponentsBuilder.fromUri(uri).build()
        .getQueryParams().getFirst("access_token");
if (token == null || token.isBlank()) {
    return null;   // handshake refuse
}

var jwt = jwtDecoder.decode(token);       // meme secret HS256 partage
var rawUserId = jwt.getClaim("user_id");

// Registre des sessions : un utilisateur peut avoir PLUSIEURS onglets ouverts.
// On garde donc un ENSEMBLE de sessions par utilisateur, et une notification
// est poussee sur toutes ses sessions a la fois.
sessionsByUser
    .computeIfAbsent(userId, ignored -> ConcurrentHashMap.newKeySet())
    .add(session);
\end{lstlisting}

\subsection{Points d'entrée exposés}

\begin{table}[H]
\centering
\small
\begin{tabularx}{\linewidth}{@{}p{1.5cm}p{6.4cm}X@{}}
\toprule
\textbf{Méth.} & \textbf{Chemin} & \textbf{Rôle} \\
\midrule
GET & \code{/api/notifications?page=0\&size=20} & Historique paginé \\
GET & \code{/api/notifications/unread-count} & Nombre de non-lues \\
PATCH & \code{/api/notifications/\{id\}/read} & Marquer comme lue \\
PATCH & \code{/api/notifications/read-all} & Tout marquer comme lu \\
DELETE & \code{/api/notifications/\{id\}} & Suppression logique \\
\emph{WS} & \code{/ws/notifications?access\_token=…} & Flux temps réel \\
\bottomrule
\end{tabularx}
\caption{Points d'entrée du Notification-Service}
\end{table}

% =============================================================================
\chapter{Fonctionnalités au-delà du cahier des charges}
\label{chap:extras}
% =============================================================================

Ce chapitre présente les quatre services supplémentaires et les mécanismes
transverses que nous avons développés sans qu'ils soient demandés. Ils ne font
pas partie du périmètre noté ; nous les présentons parce qu'ils constituent, de
notre point de vue, la partie la plus aboutie du projet.

\section{Le bus d'événements RabbitMQ}

\subsection{Un exchange topic, une enveloppe unique}

Tous les événements transitent par un \emph{unique} exchange de type
\code{topic}, nommé \code{flowforge.events}. La clé de routage est le type
d'événement lui-même (\code{task.assigned}, \code{user.registered}…), ce qui
permet à un consommateur de s'abonner par motif (\code{task.*}).

Chaque message respecte la même enveloppe, quel que soit le langage du
producteur :

\begin{lstlisting}[language=json, caption={Enveloppe commune à tous les événements}]
{
  "eventId":     "a3f1c8e2-...",        // UUID : sert de cle d'idempotence
  "eventType":   "task.assigned",       // = routing key
  "occurredAt":  "2026-07-12T12:00:00Z",
  "producer":    "task-service",
  "actorUserId": 1,                     // qui a declenche l'action
  "resource":    { "type": "task", "id": 42 },
  "data":        { }                    // charge utile propre a l'evenement
}
\end{lstlisting}

\begin{pointcle}[Trois langages, un seul contrat]
Un événement publié en Python par l'\code{auth-service} est consommé en Java par
l'\code{email-service} et en Java par le \code{search-service}, tandis qu'un
événement publié en TypeScript par le \code{task-service} est consommé par les
mêmes. Aucun code n'est partagé entre ces services — seul le \emph{contrat} de
l'enveloppe l'est. C'est précisément ce que doit permettre une architecture
microservices.
\end{pointcle}

\subsection{Fiabilité : rebut, réessais, idempotence}

Chaque consommateur possède sa propre file durable, son propre exchange de rebut
(\emph{dead-letter exchange}) et sa propre file de rebut.

\begin{lstlisting}[language=Java, caption={\code{RabbitEventConfig.java} — file durable, réessais et republication en rebut}]
@Bean
Queue searchEventsQueue(EventConsumerProperties properties) {
    return QueueBuilder.durable(properties.queue())
            // Si le message est definitivement rejete, RabbitMQ le route
            // AUTOMATIQUEMENT vers cet exchange de rebut.
            .deadLetterExchange(properties.deadLetterExchange())
            .deadLetterRoutingKey(properties.deadLetterRoutingKey())
            .build();
}

@Bean
RetryOperationsInterceptor searchRetryInterceptor(RabbitTemplate rabbitTemplate,
                                                  EventConsumerProperties properties) {
    return RetryInterceptorBuilder.stateless()
            .maxAttempts(properties.maxAttempts())         // 3 tentatives
            // Backoff EXPONENTIEL : 1 000 ms, puis x2, plafonne a 10 000 ms.
            // Une panne breve (redemarrage d'Elasticsearch) se resorbe donc
            // toute seule, sans marteler le service en panne.
            .backOffOptions(properties.retryBackoffMs(), 2.0,
                            properties.retryBackoffMs() * 10)
            // Apres la derniere tentative, le message est REPUBLIE dans la DLQ
            // au lieu d'etre detruit : il reste inspectable et rejouable.
            .recoverer(new RepublishMessageRecoverer(
                    rabbitTemplate,
                    properties.deadLetterExchange(),
                    properties.deadLetterRoutingKey()))
            .build();
}

@Bean
SimpleRabbitListenerContainerFactory rabbitListenerContainerFactory(...) {
    var factory = new SimpleRabbitListenerContainerFactory();
    factory.setMessageConverter(jsonMessageConverter);
    factory.setAcknowledgeMode(AcknowledgeMode.AUTO);
    // >>> IMPORTANT <<<
    // Sans cette ligne, un message en echec serait REMIS EN FILE indefiniment
    // et bloquerait le consommateur dans une boucle infinie ("poison message").
    factory.setDefaultRequeueRejected(false);
    factory.setAdviceChain(searchRetryInterceptor);
    return factory;
}
\end{lstlisting}

\section{Recherche full-text avec Elasticsearch}
\label{sec:search}

\subsection{Indexation pilotée par les événements}

Le \code{search-service} ne possède aucune donnée : il maintient une
\emph{projection} des données des autres services, alimentée par 18 clés de
routage RabbitMQ, dans cinq index Elasticsearch.

\mermaidfig{12-recherche-elasticsearch}{Indexation événementielle et requête filtrée par les droits}{search}{0.9}

\subsection{Le filtrage des droits est poussé dans Elasticsearch}

C'est le point dont nous sommes le plus satisfaits. Une erreur classique
consiste à interroger le moteur puis à filtrer les résultats côté application :
c'est à la fois lent (on rapatrie des documents pour les jeter) et faux (la
pagination et les scores deviennent incohérents). Ici, la liste des ayants droit
voyage \emph{dans le document indexé}, et le filtre est appliqué \emph{par le
moteur}.

\begin{lstlisting}[language=Java, caption={\code{SearchQueryService.java} — requête bool avec filtre d'autorisation}]
var body = Map.of(
    "from", safePage * safeSize,
    "size", safeSize,
    "query", Map.of("bool", Map.of(

        // --- Pertinence : recherche multi-champs avec ponderation ---
        "must", List.of(Map.of(
            "multi_match", Map.of(
                "query", q,
                // Le titre pese 4x plus que la description : une correspondance
                // dans le titre remonte donc naturellement en tete.
                "fields", List.of("title^4", "username^4", "name^4",
                                  "description^2", "body^2", "biography", "tags"),
                "type", "best_fields",
                "operator", "and"
            )
        )),

        // --- Autorisation : filtre OBLIGATOIRE, applique PAR le moteur ---
        // Un utilisateur ne peut PAS obtenir dans ses resultats une ressource
        // qu'il serait incapable d'ouvrir. Le filtrage n'est jamais fait
        // apres coup : il fait partie de la requete elle-meme.
        "filter", filters
    ))
);

filters.add(Map.of("bool", Map.of(
    "should", List.of(
        Map.of("term", Map.of("viewerUserIds", userId)),  // je peux voir la tache
        Map.of("term", Map.of("memberUserIds", userId)),  // je suis membre
        Map.of("term", Map.of("ownerUserId",   userId)),  // j'en suis proprietaire
        Map.of("term", Map.of("authorUserId",  userId)),  // j'en suis l'auteur
        Map.of("term", Map.of("type",          "user")),  // annuaire public
        Map.of("term", Map.of("visibility",    "PUBLIC")) // espace public
    ),
    "minimum_should_match", 1    // au moins UNE de ces conditions
)));
\end{lstlisting}

\subsection{Reconstruction de l'index}

L'indexation étant asynchrone, l'index peut diverger (bug de consommateur, perte
de messages, ajout d'un nouveau champ). Un point d'entrée interne
\code{POST /api/search/internal/reindex}, protégé par un en-tête
\code{X-Internal-Key}, rejoue des instantanés paginés depuis les endpoints
\code{/internal/search-snapshots} de chaque service propriétaire. L'index peut
donc être reconstruit intégralement à tout moment, sans interruption de service.

\section{Service d'emailing idempotent}

L'\code{email-service} consomme \code{user.registered} et envoie un email de
bienvenue via SMTP, à partir de gabarits HTML. Sa véritable valeur est ailleurs :
c'est un \textbf{journal de livraison}.

\begin{lstlisting}[language=Java, caption={\code{EmailDeliveryService.java} — livraison « exactement une fois »}]
// Toute tentative d'envoi est PERSISTEE dans la table email_deliveries, sous
// une cle d'idempotence deterministe.
var delivery = findOrCreateDelivery(event, email);

// Si un email identique a DEJA ete envoye, on ne le renvoie pas. RabbitMQ
// pouvant redistribuer un message, sans ce garde-fou l'utilisateur recevrait
// plusieurs fois le meme email de bienvenue.
if (delivery.getStatus() == DeliveryStatus.SENT) {
    logger.debug("Skipping duplicate email delivery {}", delivery.getId());
    return;
}
if (!email.shouldSend()) {
    markPermanentFailure(delivery, email.failureReason());
    return;
}
sendWithRetry(delivery, email);

// ---------------------------------------------------------------------------
// La cle d'idempotence est un HACHE des elements qui definissent l'unicite
// d'un envoi : l'evenement source, son type, le gabarit et le destinataire.
// Deux livraisons du MEME message produisent donc la MEME cle -- et la
// contrainte d'unicite en base fait le reste.
raw = event.eventId() + ":" + event.eventType() + ":"
    + email.templateKey() + ":" + nullToEmpty(email.recipient());
return "email:" + sha256(raw);
\end{lstlisting}

\begin{pointcle}[Pourquoi les emails de réinitialisation ne passent pas par le bus]
L'email de réinitialisation de mot de passe et les invitations sont
\emph{délibérément} envoyés en direct, sans passer par RabbitMQ. Motif : ils
contiennent des \textbf{jetons à usage unique}. Publier un tel jeton sur un bus
partagé, où il serait persisté sur disque et lisible par tout consommateur
présent ou futur, reviendrait à en faire un secret partagé. Un bus d'événements
n'est pas un canal sûr pour un secret.
\end{pointcle}

\section{Espaces de travail et pièces jointes}

\subsection{Workspace-Service}

Il possède les espaces de travail, leurs membres, leurs rôles
(\code{OWNER}, \code{ADMIN}, \code{MEMBER}, \code{VIEWER}) et leurs invitations.
C'est la source de vérité des droits, dont héritent les catégories puis les
tâches. À l'inscription, il consomme \code{user.registered} et provisionne
automatiquement un espace de travail personnel.

\subsection{Attachment-Service}

Il gère les pièces jointes des tâches et des commentaires. Les \emph{octets}
sont déposés dans un stockage objet (Azure Blob) ; la base ne contient que les
\textbf{métadonnées} et la clé de l'objet. Le téléchargement se fait par une
\textbf{URL signée à durée limitée} (15~minutes) : le fichier n'est jamais servi
par le service lui-même, qui ne devient donc jamais un goulot d'étranglement
réseau.

\mermaidfig{11-modele-workspace-notification}{Modèle de données des services Workspace, Notification, Attachment et Email}{modele-ws}{1.0}

\section{Bilan des apports hors périmètre}

\begin{table}[H]
\centering
\small
\begin{tabularx}{\linewidth}{@{}p{4.3cm}X@{}}
\toprule
\textbf{Apport} & \textbf{Intérêt technique} \\
\midrule
Bus RabbitMQ & Découplage réel : ajouter un consommateur ne modifie aucun producteur. \\
\addlinespace
Files de rebut + réessais & Aucun message n'est perdu ni rejoué en boucle. \\
\addlinespace
Idempotence & Tolérance native à la redistribution « au moins une fois ». \\
\addlinespace
Elasticsearch & Recherche full-text pondérée, avec droits appliqués dans le moteur. \\
\addlinespace
WebSocket & Notifications poussées, sans sondage côté client. \\
\addlinespace
Stockage objet & La base ne grossit pas avec les fichiers ; URLs signées. \\
\addlinespace
Publication transactionnelle & Aucun événement émis pour une écriture annulée. \\
\bottomrule
\end{tabularx}
\caption{Synthèse des fonctionnalités non demandées}
\end{table}

% =============================================================================
\chapter{Captures d'écran du frontend}
% =============================================================================

Le frontend est une \emph{single-page application} Vue~3 + TypeScript, construite
avec Vite et servie depuis la racine \code{/} par la passerelle. Elle communique
exclusivement avec \code{http://localhost:8080} : aucune URL de service n'est
codée en dur.

\section{Authentification}

\capture{01-login}{\textbf{Page de connexion.} L'utilisateur saisit son
identifiant et son mot de passe. En cas d'échec, l'erreur renvoyée par
l'\code{auth-service} s'affiche en ligne sous le formulaire. Si le compte a la
double authentification activée, le serveur répond \code{202} avec
\code{two\_factor\_required} et l'interface demande alors le code TOTP. Les
jetons reçus (\code{access} et \code{refresh}) sont stockés côté navigateur,
puis injectés automatiquement dans chaque appel ultérieur.}{login}

\capture{02-register}{\textbf{Page d'inscription.} Trois champs seulement :
nom d'utilisateur, email et mot de passe (8 caractères minimum, vérifié côté
serveur). À la validation, le compte est créé, l'utilisateur est connecté
automatiquement, puis redirigé vers l'écran de préparation de l'espace de
travail.}{register}

\capture{03-preparing-workspace}{\textbf{Préparation de l'espace de travail.}
Cet écran est la \emph{preuve visible} de la chaîne de provisionnement
asynchrone. Après l'inscription, le profil et l'espace personnel n'existent pas
encore : ils seront créés par deux consommateurs RabbitMQ distincts réagissant à
l'événement \code{user.registered}. L'interface sonde donc le backend toutes les
1,2~seconde jusqu'à ce que les deux ressources apparaissent, en affichant l'état
d'avancement, puis bascule vers l'accueil.}{preparing}

\capture{04-forgot-password}{\textbf{Réinitialisation du mot de passe.}
Déroulement en deux temps sur le même écran : saisie de l'email (le serveur
répond toujours par un message neutre, afin de ne pas révéler si l'adresse
existe), puis saisie du jeton reçu et du nouveau mot de passe.}{forgot}

\section{Tableau de bord et tâches}

\capture{05-home}{\textbf{Accueil.} Vue de synthèse de l'espace de travail
courant : tâches assignées à l'utilisateur, catégories récentes avec leur
avancement, échéances proches, et membres de l'espace. La barre supérieure
contient le sélecteur d'espace, la recherche globale, le bouton de création et la
cloche de notifications avec son compteur de non-lues.}{home}

\capture{06-your-work}{\textbf{« Votre travail ».} Liste de toutes les tâches
assignées à l'utilisateur courant, tous espaces confondus, filtrables par texte
et par statut.}{yourwork}

\capture{07-task-create}{\textbf{Création d'une tâche.} La fenêtre modale est
accessible depuis n'importe quel écran. Elle permet de renseigner le titre, la
description, la catégorie, la priorité, l'assigné et la date d'échéance. À la
validation, le \code{task-service} publie \code{task.created}, ce qui déclenche
l'indexation Elasticsearch en arrière-plan.}{taskcreate}

\capture{08-task-details}{\textbf{Détail d'une tâche.} Description, sous-tâches
cochables, pièces jointes, commentaires et journal d'activité. Le panneau latéral
porte le statut, la priorité, l'assigné, l'échéance et les étiquettes. Un
utilisateur sans droit de contribution voit les contrôles désactivés : les droits
sont chargés depuis \code{GET /api/tasks/\{id\}/permissions} et appliqués à
l'interface.}{taskdetails}

\capture{09-task-comments}{\textbf{Commentaires d'une tâche.} Le fil de
discussion est servi par le \code{comment-service}, qui a préalablement vérifié
auprès du \code{task-service} que l'utilisateur a le droit de voir la tâche.
Chaque commentaire publié est indexé dans Elasticsearch et devient donc
immédiatement cherchable.}{comments}

\section{Catégories}

\capture{10-categories-list}{\textbf{Liste des catégories.} Chaque ligne affiche
la catégorie, son responsable, sa barre d'avancement (tâches terminées / total)
et sa date de mise à jour.}{catlist}

\capture{11-category-new}{\textbf{Création d'une catégorie.} Nom, clé courte
(utilisée dans les identifiants de tâches, par exemple \code{WEB-42}),
description et couleur.}{catnew}

\capture{12-category-board}{\textbf{Vue kanban d'une catégorie.} Les tâches sont
réparties en cinq colonnes selon leur statut. Le changement de statut est
immédiatement persisté et journalisé dans \code{task\_activities}. Un
utilisateur en lecture seule voit le sélecteur désactivé.}{board}

\capture{13-category-overview}{\textbf{Synthèse d'une catégorie.} Compteurs
(total, en cours, terminées, avec échéance) et informations générales.}{catoverview}

\capture{14-category-members}{\textbf{Membres et invitations d'une catégorie.}
L'invitation génère un jeton à usage unique, affiché \emph{une seule fois} :
seul son haché SHA-256 est conservé en base.}{catmembers}

\section{Statistiques}

\capture{15-statistics}{\textbf{Tableau de bord des statistiques.} Servi par le
\code{statistics-service}, qui agrège à la volée les tâches visibles par
l'utilisateur. On y trouve le taux de complétion, la répartition par statut, par
priorité et par catégorie, ainsi que la tendance de complétion des derniers
jours. Les filtres en haut de page sont recombinables et relancent l'agrégation.}{stats}

\section{Recherche}

\capture{16-search}{\textbf{Recherche globale.} La saisie est débattue de
250~ms, puis interroge Elasticsearch via \code{/api/search}. Les résultats sont
regroupés par type (tâche, catégorie, espace, utilisateur, commentaire). Le
filtre d'autorisation étant appliqué \emph{dans} le moteur, aucun résultat
inaccessible ne peut apparaître.}{search}

\section{Notifications}

\capture{17-notifications}{\textbf{Centre de notifications.} Historique paginé,
servi par le \code{notification-service} depuis MySQL. Actions disponibles :
marquer comme lu, tout marquer comme lu, supprimer.}{notifs}

\capture{18-notification-realtime}{\textbf{Notification temps réel.} Lorsqu'un
collègue assigne une tâche, l'événement \code{task.assigned} est consommé, la
notification est persistée puis poussée sur la WebSocket ouverte. Le badge de la
cloche s'incrémente \emph{sans rechargement de page ni sondage}.}{realtime}

\section{Profil et sécurité}

\capture{19-profile}{\textbf{Profil utilisateur.} Photo (JPG/PNG/WebP, 2~Mo
maximum, validée côté serveur puis déposée dans le stockage objet), biographie,
fuseau horaire et langue.}{profile}

\capture{20-security}{\textbf{Sécurité.} Changement de mot de passe : à la
validation, \emph{toutes} les sessions sont révoquées et l'utilisateur est
redirigé vers la connexion. La double authentification se configure depuis cet
écran.}{security}

\capture{21-sessions}{\textbf{Sessions actives.} Chaque jeton de rafraîchissement
émis apparaît ici, avec son appareil, son adresse IP et sa date. Chaque session
peut être révoquée individuellement — ce que permet la matérialisation des
sessions en base (voir \code{AuthSession}).}{sessions}

\section{Espaces de travail}

\capture{22-workspace-settings}{\textbf{Paramètres d'un espace de travail.}
Écran réservé aux rôles \code{OWNER} et \code{ADMIN} : visibilité, gestion des
membres, changement de rôle, invitations, archivage.}{wssettings}

\section{Infrastructure}

\capture{23-rabbitmq}{\textbf{Console d'administration RabbitMQ}
(\code{localhost:15672}). On y voit l'exchange \code{flowforge.events}, les cinq
files de consommateurs et leurs files de rebut associées, ainsi que le débit de
messages.}{rabbitmq}

\capture{24-elasticsearch}{\textbf{Index Elasticsearch} (\code{localhost:9200}).
Les cinq index (\code{flowforge\_tasks}, \code{flowforge\_comments},
\code{flowforge\_categories}, \code{flowforge\_workspaces},
\code{flowforge\_users}) et le nombre de documents indexés.}{elastic}

\capture{25-welcome-email}{\textbf{Email de bienvenue.} Envoyé par
l'\code{email-service} en réaction à \code{user.registered}, à partir d'un
gabarit HTML. Chaque envoi est journalisé dans la table \code{email\_deliveries}
sous une clé d'idempotence, ce qui garantit qu'un même email n'est jamais envoyé
deux fois.}{email}

\capture{26-docker-ps}{\textbf{Conteneurs en cours d'exécution}
(\code{docker compose ps}). Les 19 conteneurs de la plateforme : 11 services,
9 bases, RabbitMQ, Elasticsearch, la passerelle et le frontend.}{docker}

% =============================================================================
\appendix
\chapter{Annexes}
% =============================================================================

\section{Fichier \texttt{docker-compose.yml}}
\label{ann:compose}

Ce fichier unique orchestre l'intégralité de la plateforme. La commande
\code{docker compose up -d --build} suffit à démarrer les 19 conteneurs.

\lstinputlisting[language=yaml, basicstyle=\ttfamily\tiny, numbers=left]{annexes/docker-compose.yml}

\section{Fichier \texttt{nginx.conf}}
\label{ann:nginx}

Configuration de la passerelle API : service de la SPA depuis la racine et
\emph{reverse proxy} de chaque microservice backend.

\lstinputlisting[language=nginxconf, basicstyle=\ttfamily\scriptsize, numbers=left]{annexes/nginx.conf}

\section{Code source des diagrammes (Mermaid)}
\label{ann:mermaid}

Tous les diagrammes de ce rapport ont été écrits en \textbf{Mermaid}, puis
rendus en PDF vectoriel par \code{mermaid-cli} (script
\code{build-diagrams.sh}). Les sources sont reproduites ci-dessous.

\subsection{Architecture globale (figure~\ref{fig:archi})}
\lstinputlisting[numbers=none, basicstyle=\ttfamily\tiny]{diagrams/01-architecture-globale.mmd}

\subsection{Routage NGINX (figure~\ref{fig:routage})}
\lstinputlisting[numbers=none, basicstyle=\ttfamily\tiny]{diagrams/02-routage-nginx.mmd}

\subsection{Topologie Docker (figure~\ref{fig:docker})}
\lstinputlisting[numbers=none, basicstyle=\ttfamily\tiny]{diagrams/03-topologie-docker.mmd}

\subsection{Séquence d'inscription (figure~\ref{fig:seq-inscription})}
\lstinputlisting[numbers=none, basicstyle=\ttfamily\tiny]{diagrams/04-sequence-inscription.mmd}

\subsection{Cycle de vie d'une tâche (figure~\ref{fig:seq-tache})}
\lstinputlisting[numbers=none, basicstyle=\ttfamily\tiny]{diagrams/05-sequence-tache.mmd}

\subsection{Bus d'événements (figure~\ref{fig:bus})}
\lstinputlisting[numbers=none, basicstyle=\ttfamily\tiny]{diagrams/06-bus-evenements.mmd}

\subsection{Détection des échéances (figure~\ref{fig:deadline})}
\lstinputlisting[numbers=none, basicstyle=\ttfamily\tiny]{diagrams/07-sequence-deadline.mmd}

\subsection{Modèle Auth \& User (figure~\ref{fig:modele-auth})}
\lstinputlisting[numbers=none, basicstyle=\ttfamily\tiny]{diagrams/08-modele-auth-user.mmd}

\subsection{Modèle Task (figure~\ref{fig:modele-task})}
\lstinputlisting[numbers=none, basicstyle=\ttfamily\tiny]{diagrams/09-modele-task.mmd}

\subsection{Modèle Category \& Comment (figure~\ref{fig:modele-cat})}
\lstinputlisting[numbers=none, basicstyle=\ttfamily\tiny]{diagrams/10-modele-category-comment.mmd}

\subsection{Modèle Workspace, Notification, Attachment, Email (figure~\ref{fig:modele-ws})}
\lstinputlisting[numbers=none, basicstyle=\ttfamily\tiny]{diagrams/11-modele-workspace-notification.mmd}

\subsection{Recherche Elasticsearch (figure~\ref{fig:search})}
\lstinputlisting[numbers=none, basicstyle=\ttfamily\tiny]{diagrams/12-recherche-elasticsearch.mmd}

\subsection{Flux des statistiques (figure~\ref{fig:stats})}
\lstinputlisting[numbers=none, basicstyle=\ttfamily\tiny]{diagrams/13-flux-statistiques.mmd}

\section{Démarrage de la plateforme}

\begin{lstlisting}[language=bash, numbers=none]
# 1. Cloner le depot
git clone <url-du-depot> && cd microservices-PFE

# 2. Demarrer les 19 conteneurs (premiere execution : ~5 min de build)
docker compose up -d --build

# 3. Verifier que tout tourne
docker compose ps

# 4. Ouvrir l'application -- UNE SEULE adresse a connaitre
#    http://localhost:8080

# Interfaces d'administration (optionnelles) :
#    http://localhost:15672   console RabbitMQ (flowforge / flowforge_password)
#    http://localhost:9200    API Elasticsearch
\end{lstlisting}

\end{document}
